\chapter{三角函数}

\section{正弦函数的图像与性质（1）}
\zsd 正弦函数的定义、图象与性质 \\
A组：
\begin{enumerate}
    \item 已知正弦函数过点$\left(\dfrac{\pi}{6}, m\right)$，则$m = $\blkda{$\dfrac{1}{2}$}
    \item 函数$y=\sin{x}$的单调增区间为\blkda[4]{$\left[2k\pi - \dfrac{\pi}{2}, 2k\pi+\dfrac{\pi}{2}\right], k \in \mathbf{Z}$}
    \item 函数$y = sin2x$的严格减区间为\blkda[4]{$\left[k\pi+\dfrac{\pi}{4}, k\pi+\dfrac{3\pi}{4}\right], k \in \mathbf{Z}$}
    \item 函数$y = sinx, x \in \left[-\dfrac{\pi}{6}, \dfrac{2\pi}{3}\right]$的值域是\blkda{$\left[-\dfrac{1}{2}, 1\right]$}
    \item 函数$y = sin\left(x - \dfrac{\pi}{4}\right)$的图象的对称轴为\blkda[4]{直线$x = k\pi+\dfrac{3\pi}{4}, k \in \mathbf{Z}$}，对称中心为\blkda[4]{点$\left(k\pi + \dfrac{\pi}{4}, 0\right), k\in \mathbf{Z}$}
\end{enumerate}
B组:
\begin{enumerate}
\item 若函数$y = 3sin(2x+\theta)$为偶函数，则$\theta$的取值集合是\blkda[4]{$\{x | x = k\pi + \dfrac{\pi}{2}, k \in \mathbf{Z} \}$}
\begin{jieda}
    $y = 3sin(2x+\theta) = 3cos(2x+\theta-\dfrac{\pi}{2})$，故$\theta - \dfrac{\pi}{2} = k\pi, k\in \mathbf{Z}$，故$\theta \in \{x | x = k\pi + \dfrac{\pi}{2}, k \in \mathbf{Z} \}$
\end{jieda}
\item 函数$y = |sinx|$的周期为\blkda[3]{$k\pi(k \in \mathbf{Z}, k \neq 0)$}
\item 函数$f(x)=\sin{x}+2\abs{\sin{x}}$, $x\in[0,2\pi]$的图像与直线$y=k$有且仅有两个不同的交点, 则$k$的取值范围是\blkda{$(1,3)$}. 

\item 方程$\lg{x}-\sin{x}=0$的根个数是\blkda{3}. 



\item 设函数$f(x)=\sin3x+\abs{\sin3x}$, 则$f(x)$为\dotfill(\kh{A})
\xx
{周期函数, 最小正周期为$\dfrac{2\pi}{3}$}
{周期函数, 最小正周期为$\dfrac{\pi}{3}$}
{周期函数, 数小正周期为$2\pi$}
{非周期函数}

\item 已知函数$y=\sin\left(x-\dfrac{\pi}{12}\right)\cos\left(x-\dfrac{\pi}{12}\right)$, 则下列判断正确的是\dotfill(\kh{B})
\xx
{此函数的最小正周期为$2\pi$, 其图像的一个对称中心是$\left(\dfrac{\pi}{12},0\right)$}
{此函数的最小正周期为$\pi$, 其图像的一个对称中心是$\left(\dfrac{\pi}{12},0\right)$}
{此函数的最小正周期为$2\pi$, 其图像的一个对称中心是$\left(\dfrac{\pi}{6},0\right)$}
{此函数的最小正周期为$\pi$, 其图像的一个对称中心是$\left(\dfrac{\pi}{6},0\right)$}
\begin{jieda}
$y=\frac{1}{2}\sin(2x-\frac{\pi}{6})$. 
\end{jieda}
\item 若在$\left[0,\frac{\pi}{2}\right]$内有两个不同的实数值满足等式$\cos2x+\sqrt{3}\sin2x=k+1$, 则$k$的取值范围是\dotfill(\kh{B})
\xx
{$0<k\leq 1$}
{$0\leq k<1$}
{$-3\leq k\leq 1$}
{$k\leq 1$}
\begin{jieda}
$\iff \sin(2x+\frac{\pi}{6})=\frac{k+1}{2}$. \\
$x\in[0,\frac{\pi}{2}]\iff 2x+\frac{\pi}{6}\in[\frac{\pi}{6},\frac{7\pi}{6}]$. \\
考虑到两解, 故$2x+\frac{\pi}{6}\in[\frac{\pi}{6},\frac{\pi}{2})\cup(\frac{\pi}{2},\frac{5\pi}{6}]$, 于是$\frac{k+1}{2}\in[\frac{1}{2},1)$. 
\end{jieda}
\item 画出下列函数的图像，并求函数的单调区间: 
\begin{enumerate}[(1)]
\item $y=\sin(x+\tfrac{\pi}{3})$;
\item $y=\sin{x}-\cos{x}$. 
\end{enumerate}
\begin{jieda}
\begin{enumerate}[(1)]
\item 单调递增: $x+\frac{\pi}{3}\in[2k\pi-\frac{\pi}{2},2k\pi+\frac{\pi}{2}]\iff x\in[2k\pi-\frac{5\pi}{6},2k\pi+\frac{\pi}{6}](k\zz)$. \\
单调递减: $[2k\pi+\frac{\pi}{6},2k\pi+\frac{7\pi}{6}](k\zz)$. 
\item $y=\sqrt{2}\sin(x-\frac{\pi}{4})$, 易得单调递增: $[2k\pi-\frac{\pi}{4},2k\pi+\frac{3\pi}{4}](k\zz)$, 单调递减: $[2k\pi+\frac{3\pi}{4},2k\pi+\frac{7\pi}{4}](k\zz)$. 
\end{enumerate}
\end{jieda}

\item 已知函数$f(x) = sin\omega x+acos\omega x (\omega > 0)$的最小正周期为$2\pi$
\begin{enumerate}
    \item 求$\omega$的值
    \item 已知直线$x = -\dfrac{\pi}{4}$是函数$f(x)$的一条对称轴，求$f(x)$的最大值与最小值
\end{enumerate}
\begin{jieda}
    $f(x) = \sqrt{a^2+1}sin(\omega x + \varphi)$（$cos\varphi = \dfrac{1}{\sqrt{a^2+1}}, sin\varphi = \dfrac{a}{\sqrt{a^2+1}}$）$\dfrac{2\pi}{\omega} = T = 2\pi$，故$\omega = 1$
    \begin{enumerate}
        \item 方法1 \\
         $f(x) = sinx + a \cdot cosx$，因为$x = -\dfrac{\pi}{4}$是函数的一条对称轴，因此$f(-\dfrac{\pi}{4}) = \pm \sqrt{a^2+1}$，即$-\dfrac{\sqrt{2}}{2} + a \cdot \dfrac{\sqrt{2}}{2} = \pm \sqrt{a^2+1}$，即$a-1 = \pm \sqrt{2(a^2+1)}$，即$a^2-2a+1 = 2(a^2+1)$，即$a^2+2a+1 = 0$，因此$a = -1$，故$f(x)$的最大值为$\sqrt{2}$，最小值为$-\sqrt{2}$
        \item 方法2 \\
        $f(x) = \sqrt{a^2+1} sin(x+\varphi)$，因为$x = -\dfrac{\pi}{4}$是函数的一条对称轴，因此$-\dfrac{\pi}{4} + \varphi = k\pi + \dfrac{\pi}{2}, k \in \mathbf{Z}$，因此$\varphi = k\pi + \dfrac{3\pi}{4}, k \in \mathbf{Z}$因此$tan\varphi = a = -1$，因此$f(x) = \sqrt{2}sin(x - \dfrac{\pi}{4})$，故其最大值为$\sqrt{2}$，最小值为$-\sqrt{2}$
        \item 方法3 \\
        $f(x) = sinx + a \cdot cosx$，因为$x = -\dfrac{\pi}{4}$是函数的一条对称轴，因此$f(0) = f(-\dfrac{\pi}{2})$，故$a = -1$，经检验$f(x) = sinx - cosx$符合要求，故$f(x)$的最大值为$\sqrt{2}$，最小值为$-\sqrt{2}$
    \end{enumerate}
    当$x = 2k\pi - \dfrac{\pi}{4}, k \in \mathbf{Z}$的时候取得最小值，当$x = 2k\pi + \dfrac{\pi}{4}, k \in \mathbf{Z}$的时候取得最大值
\end{jieda}

\item 已知函数$f(x)=\cos^2x-2\sin{x}\cos{x}-\sin^2x$. 
\begin{enumerate}[(1)]
\item 求$f(x)$的最小正周期;
\item 求$f(x)$的最大值和最小值及对应的$x$的值.
\end{enumerate}
\begin{jieda}
$f(x)=\sqrt{2}\cos(2x+\frac{\pi}{4})=\sqrt{2}\sin(2x+\frac{3\pi}{4})$. 
\begin{enumerate}[(1)]
\item $T=\pi$; 
\item $\max=f(k\pi-\frac{\pi}{8})=\sqrt{2}$, $\min=f(k\pi-\frac{5\pi}{8})=-\sqrt{2}$. 
\end{enumerate}
\end{jieda}


\item 已知函数$f(x)=2\sin\frac{x}{4}\cos\frac{x}{4}-2\sqrt{3}\sin^2\frac{x}{4}+\sqrt{3}$. 
\begin{enumerate}[(1)]
\item 求函数$f(x)$的最小正周期及最值;
\item 令$g(x)=f\left(x+\frac{\pi}{3}\right)$, 判断函数$g(x)$的奇偶性, 并说明理由. 
\end{enumerate}
\begin{jieda}
$f(x)=2\sin(\frac{x}{2}+\frac{\pi}{3})$. 
\begin{enumerate}
\item $T=4\pi$, $\min=f(4k\pi-\frac{5\pi}{3})=-2$, $\max=f(4k\pi+\frac{\pi}{3})=2$; 
\item $g(x)=2\cos\frac{x}{2}$, 偶函数. 
\end{enumerate}
\end{jieda}
\end{enumerate}


\section{正弦函数的图像与性质（2）}
\zsd 正弦函数的定义、图象与性质 \\
A组：
\begin{enumerate}
    \item 函数$y = sin\left(2x+\dfrac{\pi}{3}\right)$的图象的对称轴方程是\blkda[3.5]{$x = \dfrac{k\pi}{2}+\dfrac{\pi}{12}, k \in \mathbf{Z}$}，对称中心的坐标是\blkda[3.5]{$\left(\dfrac{k\pi}{2}-\dfrac{\pi}{6}, 0\right), k \in \mathbf{Z}$}
    \item 函数$y = 2sin\left(2x-\dfrac{\pi}{3}\right)$的单调递增区间为\blkda[4]{$\left[k\pi - \dfrac{\pi}{12}, k\pi + \dfrac{5\pi}{12}\right], k \in \mathbf{Z}$}
    \item 函数$y = sin\left(2x+\dfrac{5\pi}{2}\right)$的图象的对称轴方程是\blkda[3.5]{$x = \dfrac{k\pi}{2}-\pi, k \in \mathbf{Z}$}
    \item 函数$y = 2 - 3sinx$，当$x = $\blkda[4]{$2k\pi + \dfrac{\pi}{2}, k \in \mathbf{Z}$}，$y$取最小值\blkda{$-1$}; \\
    当$x = $\blkda[3]{$2k\pi + \dfrac{3\pi}{2}, k \in \mathbf{Z}$}时，$y$取最大值\blkda{5}
    \item 函数$y = \sqrt{sinx}$，当$x = $\blkda[3]{$k\pi, k \in \mathbf{Z}$}时，有最小值\blkda{0}
    \item 设$x \in \mathbf{Z}$，则函数$y = sin\dfrac{\pi x}{3}$的值域为\blkda[3]{$\{0, \dfrac{\sqrt{3}}{2}, -\dfrac{\sqrt{3}}{2}\}$}
    \item 函数$y = \dfrac{sinx - 1}{sinx + 2}$的值域是\blkda{$[-2, 0]$}
    \item 函数$y = sin^2x + sinx + 1$的值域为\blkda{$\left[\dfrac{3}{4}, 3\right]$}
    \item 函数$f(x) = sin^2x + 2sinx cosx + 3cos^2x$的值域是\blkda[3]{$[2-\sqrt{2}, 2+\sqrt{2}]$}
    \item 判断下列函数的奇偶性，并写出最小正周期
    \begin{dingautolist}{172}
        \item $y = sin\dfrac{3}{2}x, x \in \mathbf{R}$；\blkda[3]{奇函数}，\blkda{$\dfrac{4\pi}{3}$}
        \item $y = sin\left(\dfrac{1}{2}x - \dfrac{\pi}{3}\right), x \in \mathbf{R}$；\blkda[3]{非奇非偶函数}，\blkda{$4\pi$}
        \item $f(x) = \pi sin\pi x$；\blkda[3]{奇函数}，\blkda{2}
        \item $y = cos2x + sinx \cdot cosx$；\blkda[3]{非奇非偶函数}，\blkda{$\pi$}
        \item $f(x) = sin^2x + 2sinx\cdot cosx + 3cos^2x$；\blkda[3]{非奇非偶函数}，\blkda{$\pi$}
    \end{dingautolist}
\end{enumerate}
B组：
\begin{enumerate}
    \item 函数$y = log_{\frac{1}{2}}sin\left(\dfrac{x}{2} - \dfrac{\pi}{3}\right)$的单调递增区间为\blkda[6]{$\left[\dfrac{5\pi}{3}+4k\pi, \dfrac{8\pi}{3}+4k\pi \right), k \in \mathbf{Z}$}
    \item 若$f(x)=2\sin\omega{x}(0<\omega<1)$在区间$\left[0,\frac{\pi}{3}\right]$上的最大值为$\sqrt{2}$, 则$\omega$=\blkda{$\frac{3}{4}$}
    \item 关于函数$f(x) = 2sin\left(3x - \dfrac{3\pi}{4}\right)$，以下说法：
    \begin{dingautolist}{172}
        \item 其最小正周期为$\dfrac{2\pi}{3}$
        \item 图像关于点$\left(\dfrac{\pi}{4}, 0\right)$对称
        \item 直线$x = -\dfrac{\pi}{4}$是其一条对称轴
    \end{dingautolist}. 其中说法正确的序号是\blkda{\ding{172}, \ding{173}, \ding{174}}
    \item 判断下列函数的奇偶性，并写出最小正周期
    \begin{dingautolist}{172}
        \item $f(x) = sin^6x + cos^6x$；\blkda[3]{偶函数}，\blkda{$\dfrac{\pi}{2}$}
        \item $f(x) = asin^2x + bcos^2x(a^2+b^2 \neq 0, b \neq a)$；\blkda[3]{偶函数}，\blkda{$\pi$}
    \end{dingautolist}
    \item 函数$y=\cos2x+\sin^2x$的值域是\blkda{$[0,1]$}
    \item 函数$y=-2\sin(\frac{\pi}{4}-3x)$的单调递增区间为\blkda[5]{$[\frac{2k\pi}{3}-\frac{\pi}{12},\frac{2k\pi}{3}+\frac{\pi}{4}](k\zz)$}
    \item 函数$y=\sin{x}\cdot\abs{\cos{x}}$的最小正周期为\blkda{$2\pi$}
    \item 函数$f(x)=\abs{\sin{x}+\cos{x}}$的最小正周期是 \blkda{$\pi$}
    \item 已知函$f(x)=\sin\left(\omega{x}+\frac{\pi}{3}\right)(\omega>0)$的最小正周期为$\pi$, 则该函数的图像\dotfill(\kh{A})
    \xx
    {关于点$(\frac{\pi}{3},0)$对称}
    {关于直线$x=\frac{\pi}{4}$对称}
    {关于点$(\frac{\pi}{4},0)$对称}
    {关于直线$x=\frac{\pi}{3}$对称}
    \item 已知函数$f(x)=2\cos^2x+\sqrt{3}\sin2x+a$, 若$x\in\left[0,\frac{\pi}{2}\right]$, 且$\abs{f(x)}<2$, 求$a$的取值范围.
    \begin{jieda}
    $f(x)=2\sin(2x+\frac{\pi}{6})+a+1$. \\
    $x\in[0,\frac{\pi}{2}]$, 则$2x+\frac{\pi}{6}\in[\frac{\pi}{6},\frac{7\pi}{6}]$, 即$f(x)\in[a,a+3]$. \\
    于是$-2<a<a+3<2$, 故$a\in(-2,-1)$. 
    \end{jieda}
    \item 作出下列函数在给定区间上的图像: 
    \begin{enumerate}[(1)]
    \item $y=\sin{x}+\abs{\sin{x}}(x\in[-\pi,\pi])$; 
    \item $y=\abs{\sin{x}}+\abs{\cos{x}}(x\in[-\pi,\pi])$.
    \end{enumerate}
    \begin{jieda}
    五点作图法未必一直有效. 逐个象限分段考虑.
    \end{jieda}
    \item 已知函数$f(x)=\sin^2\omega{x}+\sqrt{3}\sin\omega{x}\sin\left(\omega{x}+\frac{\pi}{2}\right)(\omega >0)$的最小正周期为$\pi$.
    \begin{enumerate}[(1)]
    \item 求$\omega$的值;
    \item 求函数$f(x)$在区间$\left[0,\frac{2\pi}{3}\right]$上的取值范围.
    \end{enumerate}
    \begin{jieda}
    $f(x)=\sin(2\omega{x}-\frac{\pi}{6})+\frac{1}{2}$. 
    \begin{enumerate}[(1)]
    \item $\omega=1$; 
    \item $[0,\frac{3}{2}]$. 
    \end{enumerate}
    \end{jieda}
    \item 已知函数$f(x)=a(2\cos^2\frac{x}{2}+\sin{x})+b$. 
    \begin{enumerate}[(1)]
    \item 当$a=1$时, 求$f(x)$的单调递增区间;
    \item 当$a<0$时, 且$x\in[0,\pi]$时, $f(x)$的值域是$[3,4]$, 求$a,b$的值.
    \end{enumerate}
    \begin{jieda}
    $f(x)=\sqrt{2}a\sin(x+\frac{\pi}{4})+a+b$. 
    \begin{enumerate}
    \item $[2k\pi-\frac{3\pi}{4},2k\pi+\frac{\pi}{4}](k\zz)$. 
    \item $x\in[0,\pi]$, 则$x+\frac{\pi}{4}\in[\frac{\pi}{4},\frac{5\pi}{4}]$, 故$\sqrt{2}\sin(x+\frac{\pi}{4})\in[-1,\sqrt{2}]$. \\
    因为$a<0$, 故$f(x)\in[\sqrt{2}a+a+b,b]$, 故$\sqrt{2}a+a+b=3\land b=4\land a<0$, \\
    解得$a=1-\sqrt{2},b=4$. 
    \end{enumerate}
    \end{jieda}
\end{enumerate}
C组：
\begin{enumerate}
    \item 已知函数$f(x)=\dfrac{\sin(\pi x)-\cos(\pi x)+2}{\sqrt{x}}\left(\dfrac{1}{4}\leq x\leq \dfrac{5}{4}\right)$, 则$f(x)$的最小值为\blkda{$\frac{4\sqrt{5}}{5}$}. 
    \begin{jieda}
    $f(x)=\dfrac{\sqrt{2}\sin(\pi{x}-\frac{\pi}{4})+2}{\sqrt{x}}\geq\dfrac{2}{\sqrt{x}}\geq\dfrac{2}{\sqrt{\frac{5}{4}}}=\dfrac{4}{\sqrt{5}}$. 
    等号成立的条件$x=\frac{5}{4}$. 
    \end{jieda}
    \item 若函数$f(x)=2a\sin(2x+\varphi)\left(0<\varphi<2\pi\right)$的值域为$[-2,2]$, 且在区间$\left[-\frac{5\pi}{12},\frac{\pi}{12}\right]$上是减函数, 
    求常数$a$和$\varphi$的值. 
    \begin{jieda}
    显然$a\neq 0$(否则值域为$\{0\}$), 则$T=\pi$. \\
    而$\frac{\pi}{12}-(-\frac{5\pi}{12})=\frac{\pi}{2}=\frac{T}{2}$(半个周期), 且在区间$\left[-\frac{5\pi}{12},\frac{\pi}{12}\right]$上减, \\
    故$f(-\frac{5\pi}{12})=\max\land f(\frac{\pi}{12})=\min$, 且$\abs{a}=1$. 
    \begin{enumerate}[(1)]
    \item 当$a=1$时, $\begin{cases}
    2\cdot(-\frac{5\pi}{12})+\varphi=2k\pi+\frac{\pi}{2}\\
    2\cdot\frac{\pi}{12}+\varphi=2k\pi+\frac{3\pi}{2}
    \end{cases}$, 
    则$\varphi=2k\pi+\frac{4\pi}{3}$. \\
    因为$\varphi\in(0,2\pi)$, 故$\varphi=\frac{4\pi}{3}$. 
    
    \item 当$a=-1$时, $\begin{cases}
    2\cdot(-\frac{5\pi}{12})+\varphi=2k\pi-\frac{\pi}{2}\\
    2\cdot\frac{\pi}{12}+\varphi=2k\pi+\frac{\pi}{2}
    \end{cases}$, 
    则$\varphi=2k\pi+\frac{\pi}{3}$. \\
    因为$\varphi\in(0,2\pi)$, 故$\varphi=\frac{\pi}{3}$. 
    \end{enumerate}
    
    综上, $(a,\varphi)=(1,\frac{4\pi}{3}),(-1,\frac{\pi}{3})$. 
    \end{jieda}
\end{enumerate}


\section{余弦函数的图像与性质（1）}
\zsd 余弦函数的图像与性质 \\

A组 :
\begin{enumerate}
    \item 函数$f(x)=\cos x, x \in[-\frac{\pi}{4}, \frac{5 \pi}{6}]$ 的值域为 \blkda[2]{$[-\dfrac{\sqrt{3}}{2},1]$}.

    \item 函数$f(x)=\cos{x}+\cos{2x}$的值域是\blkda[2]{$[-\dfrac{9}{8},2]$}.

    \item 函数$y=\cos{2x},x\in [-\dfrac{\pi}{6},\dfrac{2\pi}{3}]$的单调增区间是\blkda[4]{$[-\dfrac{\pi}{6},0]$和$[\dfrac{\pi}{2},\dfrac{2\pi}{3}]$}.

    \item 已知函数$y = 2\cos (2 x-\frac{\pi}{4})$  ，当函数值为 -2 时，自变量 $x$ 的取值集合为 \blkda[4]{$\{x|x=\dfrac{5\pi}{8}+k\pi,k\in Z\}$}.

    \item 函数$y=\dfrac{\sqrt{\sin{x}}}{1+\cos{x}}$的定义域是\blkda[4]{$[2k\pi,2k\pi+\pi)(k\zz)$}. 
\end{enumerate}

B组：
\begin{enumerate}
    \item 若函数$f(x)=\cos(3x+\varphi)$的图像关于原点对称, 则$\varphi$的取值集合为\blkda[4]{$\{x\mid x=k\pi+\frac{\pi}{2},k\zz\}$}. 

    \item 函数$y=\dfrac{2+\cos{x}}{2-\cos{x}}$的值域是\blkda{$[\frac{1}{3},3]$}.

    \item 函数$y=\sqrt{2}\cos(\frac{\pi}{4}-2x)-3$的单调递增区间是\blkda[4.5]{$[k\pi-\frac{3\pi}{8},k\pi+\frac{\pi}{8}](k\zz)$}. 
\begin{jieda}
$y=\sqrt{2}\cos(2x-\frac{\pi}{4})-3$, 则$2x-\frac{\pi}{4}\in[2k\pi-\pi,2k\pi]$, 故$x\in[k\pi-\frac{3\pi}{8},k\pi+\frac{\pi}{8}](k\zz)$. 
\end{jieda}

    \item 函数$f(x)=\cos^4(\omega{x}-\frac{\pi}{6})$的最小正周期为$\frac{\pi}{5}$, 其中$\omega>0$, 则$\omega$=\blkda{5}. 
\begin{jieda}
$f(x)=\left(\dfrac{1+\cos(2\omega{x}-\frac{\pi}{3})}{2}\right)^2$. \\
因为$\dfrac{1+\cos(2\omega{x}-\frac{\pi}{3})}{2}\geq 0$, 故外层的平方并没有改变周期. \\
于是$\frac{\pi}{5}=\frac{2\pi}{2\omega}$, 即$\omega=5$. 
\end{jieda}

    \item 函数$y=(\cos{x}-m)^2-1$, 当$\cos{x}=-1$时取最大值, 当$\cos{x}=m$时取最小值, 则实数$m$的取值范围是\blkda{$[0,1]$}. 
\begin{jieda}
$\cos{x}\in[-1,1]$. 作图易得. 
\end{jieda}

\item 已知函数$y=2\cos{x}(0\leq x\leq 2\pi)$的图像和直线$y=2$围成一个封闭的图形, 则其面积为\blkda{$4\pi$}.


\item 已知关于 $x$ 的不等式 $\cos ^2 x-4 \cos x+a \geq 1$ 在 $\left[0, \frac{\pi}{2}\right]$ 内恒成立，则实数 $a$ 的取值范围是\blkda[3]{$[4,+\infty)$} ．

\item 函数 $y=2 \cos ^2 x+2 \sqrt{3} \sin x \cos x+5, x \in\left[-\frac{\pi}{6}, \frac{\pi}{3}\right]$ 的最大值是 \blkda[1]{8} ．

\item 已知 $k<-4$ ，则函数 $y=\cos 2 x+k(\cos x-1)$ 的最小值是\blkda[1]{1}．

\item 如果函数 $y=3 \cos (2 x+\varphi)$ 的图像关于点 $\left(\frac{4 \pi}{3}, 0\right)$ 中心对称，那么 $|\varphi|$ 的最小值为\blkda[1]{$\dfrac{\pi}{6}$} ．

\item 已知函数 $f(x)=\cos \omega x(\omega>0)$ 在区间 $[0,2 \pi]$ 有且仅有 3 个零点，则 $\omega$ 的取值范围是 \blkda[2]{$[\dfrac{5}{4},\dfrac{7}{4})$} ．

\item 已知函数$f(x)=2\cos{x}(\sin{x}-\cos{x})+1$. 
\begin{enumerate}[(1)]
\item 求函数$f(x)$的最小正周期; 
\item 求函数$f(x)$在$[\frac{\pi}{8},\frac{3\pi}{4}]$上的最大值与最小值. 
\end{enumerate}
\begin{jieda}
$f(x)=\sqrt{2}\sin(2x-\frac{\pi}{4})=\sqrt{2}\cos(2x-\frac{3\pi}{4})$. 
\begin{enumerate}[(1)]
\item $T=\pi$. 
\item $2x-\frac{\pi}{4}\in[0,\frac{5\pi}{4}]$, 故$\sin(2x-\frac{\pi}{4})\in[-\frac{\sqrt{2}}{2},1]$. \\
于是$\min=f(\frac{3\pi}{4})=-1$, $\max=f(\frac{3\pi}{8})=\sqrt{2}$. 
\end{enumerate}
\end{jieda}

\item 已知函数$f(x)=\cos^2\left(x+\frac{\pi}{12}\right)$, $g(x)=1+\frac{1}{2}\sin2x$. 
\begin{enumerate}[(1)]
\item 求函数$f(x)$图像的一条对称轴; 
\item 设$x=x_0$是函数$y=f(x)$图像的一条对称轴, 求$g(x_0)$的值; 
\item 求函数$h(x)=f(x)+g(x)$的单调递增区间. 
\end{enumerate}
\begin{jieda}
$f(x)=\dfrac{1+\cos(2x+\frac{\pi}{6})}{2}=\dfrac{1+\sin(2x+\frac{2\pi}{3})}{2}$. 
\begin{enumerate}[(1)]
\item $2x+\frac{\pi}{6}=k\pi\iff x=\frac{k\pi}{2}-\frac{\pi}{12}$. 
\item $2x_0+\frac{\pi}{6}=k\pi\iff x_0=\frac{k\pi}{2}-\frac{\pi}{12}$, 故$g(x_0)=1+\frac{1}{2}\sin(k\pi-\frac{\pi}{6})$. \\
当$k$为奇数时, $g(x_0)=\frac{5}{4}$; 当$k$为偶数时, $g(x_0)=\frac{3}{4}$. 
\item $f(x)+g(x)=\frac{1}{2}\sin(2x+\frac{\pi}{3})+\frac{3}{2}$, 易得单调递增区间为$[k\pi-\frac{5\pi}{12},k\pi+\frac{\pi}{12}](k\zz)$. 
\end{enumerate}
\end{jieda}

\item 设关于$x$的函数$y=2\cos^2x-2a\cos x-(2a+1)$的最小值为$f(a)$. 
\begin{enumerate}[(1)]
\item 试用$a$写出$f(a)$的表达式; 
\item 试确定$f(a)=\dfrac{1}{2}$的$a$值, 并对此时的$a$求出$y$的最大值. 
\end{enumerate}
\begin{jieda}
\begin{enumerate}[(1)]
\item $y=2(\cos{x}-\frac{a}{2})^2-\frac{a^2}{2}-2a+1$, $\cos{x}\in[-1,1]$, \\
故$f(a)=\begin{cases}
1, & a<-2\\
-\frac{1}{2}a^2-2a+1, & a\in[-2,2]\\
1-4a, & a>2
\end{cases}$. 

\item 易得$a=-1$, 故$y=2(\cos{x}+\frac{1}{2})^2+\frac{1}{2}$, 于是$y_{\max}=5$. 
\end{enumerate}
\end{jieda}

\item 已知：$f(a)=\frac{\sin \left(\frac{\pi}{2}+a\right)}{\cot \left(\frac{\pi}{2}-a\right)} \cdot \frac{\cot \left(\frac{3 \pi}{2}-a\right)}{\cos \left(\frac{3 \pi}{2}+a\right)} \frac{\cos (3 \pi-a)}{\cot (-a)}$ ．
（1）化简：$f(a)$ ；
（2）求函数 $y=f(2 x)-2 f(x)$ 的最小值．
\begin{jieda}

\begin{enumerate}[(1)]
    \item $f(\alpha)=\cos{\alpha}$
    \item $y=\cos{2x}-2\cos{x}=2\cos^{2}{x}-2\cos{x}-1,\cos{x}\in[-1,1]$,从而$y_{min}=-\dfrac{3}{2}$
\end{enumerate}
\end{jieda}

\end{enumerate}

\section{余弦函数的图像与性质（2）}
\zsd 正余弦函数综合 \\

\begin{enumerate}
    \item 下列说法中不正确的是\dotfill(\kh{C})
\xx
{正弦函数，余弦函数的定义域是 $\mathbf{R}$ ，值域是 $[-1,1]$}
{余弦函数当且仅当 $x=2 k \pi(k \in \mathbf{Z})$ 时，取得最大值 1}
{余弦函数在 $\left[2 k \pi+\frac{\pi}{2}, 2 k \pi+\frac{3 \pi}{2}\right](k \in \mathbf{Z})$ 上都是严格减函数}
{余弦函数在 $[2 k \pi-\pi, 2 k \pi](k \in \mathbf{Z})$ 上都是严格增函数}

\item 已知函数 $f(x)=2 \cos \left(\frac{\pi}{4}-3 x\right), x \in\left[-\frac{\pi}{2}, \frac{\pi}{2}\right]$ ，则 $f(x)$ 的单调递增区间是\blkda[3]{$\left[-\frac{\pi}{4}, \frac{\pi}{12}\right] \cup [\frac{5\pi}{12}, \frac{\pi}{2}]$}

\item 函数$y=7-4\sin{x}\cos{x}+4\cos^2x-4\cos^4x$的值域是\blkda{$[6,10]$}. 
\begin{jieda}
$y=\sin^22x-2\sin2x+7=(\sin2x-1)^2+6$. 
\end{jieda}

\item $y=\sin^6x+\cos^6x$的值域是\blkda{$[\frac{1}{4},1]$}. 
\begin{jieda}
$y=(\sin^2x+\cos^2x)(\sin^4x-\sin^2x\cos^2x+\cos^4x)=1-3\sin^2x\cos^2x=1-\frac{3}{4}\sin^22x$. 
\end{jieda}

\item $y=1+\sin{x}+\cos{x}+\sin{x}\cos{x}$的值域是\blkda[3]{$[0,\frac{3}{2}+\sqrt{2}]$}. 
\begin{jieda}
$y=1+t+\frac{t^2-1}{2}$, 其中$t=\sin{x}+\cos{x}\in[-\sqrt{2},\sqrt{2}]$. 
\end{jieda}


\item $y=\dfrac{\sin{x}\cos{x}}{1+\sin{x}+\cos{x}}$的值域是\blkda[4.5]{$[\frac{-\sqrt{2}-1}{2},-1)\cup(-1,\frac{\sqrt{2}-1}{2}]$}. 
\begin{jieda}
$y=\frac{t-1}{2}$, 其中$t=\sin{x}+\cos{x}\in[-\sqrt{2},-1)\cup(-1,\sqrt{2}]$. 
\end{jieda}


\item 函数$y=\dfrac{\sin{x}\cos{x}}{\sin{x}-\cos{x}+1}(0<x<\pi)$的值域是\blkda{$[\frac{1-\sqrt{2}}{2},1)$}. 
\begin{jieda}
$t=\sin{x}-\cos{x}=\sqrt{2}\sin(x-\frac{\pi}{4})\in(-1,\sqrt{2}]$. \\
$y=\dfrac{\frac{1-t^2}{2}}{t+1}=\frac{1-t}{2}\in[\frac{1-\sqrt{2}}{2},1)$. 
\end{jieda}

\item 若函数$y=\sin2x+a\cos2x$的图像关于直线$x=\frac{\pi}{6}$对称, 则$a$=\blkda{$\frac{\sqrt{3}}{3}$}. 
\begin{jieda}
$f(\frac{\pi}{6})=\pm\sqrt{1+a^2}$. 
\end{jieda}

\item 设$f(x)=\cos(\sin x)$与$g(x)=\sin(\cos x)$, 以下结论中, 正确的结论序号是\blkda{\ding{172}\ding{173}\ding{175}}: 
\begin{dingautolist}{172}
\item $f(x)$与$g(x)$都是偶函数; 
\item $f(x)$与$g(x)$都是周期函数; 
\item $f(x)$与$g(x)$的定义域都是$[-1,1]$; 
\item $f(x)$的值域是$[\cos1,1]$, $g(x)$的值域是$[-\sin1,\sin1]$. 
\end{dingautolist}
\begin{jieda}
一般情形下, $2\pi$是$f(\sin{x})$的周期, 也是$f(\cos{x})$的周期. 是否最小需要进一步的分析. \\
其二, 原来的对称轴还是对称轴, 至于是否增加新的对称轴, 需要进一步的分析. 对称中心的情况与$f(x)$有关. 
\end{jieda}


\item 给出函数$f(x)=\begin{cases}
\sin{x}, & \text{当}\sin{x}\geq\cos{x}\text{时}\\
\cos{x}, & \text{当}\sin{x}<\cos{x}\text{时}
\end{cases}$, 给出下列四个结论, 其中正确的结论序号是\blkda{\ding{174}}: 
\begin{dingautolist}{172}
\item 该函数的值域为$[-1,1]$; 
\item $f(x)$是周期函数, 最小正周期为$\pi$; 
\item 当且仅当$2k\pi+\pi<x<2k\pi+\frac{3\pi}{2}(k\zz)$时, $f(x)<0$; 
\item 当且仅当$x=2k\pi+\frac{\pi}{2}(k\zz)$时, 该函数取得最大值. 
\end{dingautolist}


\item 已知函数 $y=\sin x+a \cos x$ 的图象关于 $x=\frac{5}{3} \pi$ 对称，则函数 $y=a \sin x+\cos x$的图象的一条对称轴是\dotfill(\kh{C})
\xx
{$x=\frac{1}{3} \pi$}
{$x=\frac{2}{3} \pi$}
{$x=\frac{11}{6} \pi$}
{$x=\pi$}
\begin{jieda}
    设$f(x) = \sin x+a \cos x$，则$f\left(\dfrac{\pi}{2} - x\right) = a \sin x+\cos x$，因为$x=\frac{5}{3} \pi$是$f(x)$的对称轴，所以$x = -\dfrac{7\pi}{6}$是$f\left(\dfrac{\pi}{2} - x\right)$的对称轴，又因为$f(x), f\left(\dfrac{\pi}{2} - x\right) $的周期都是$2\pi$，所以$x = k\pi - \dfrac{7\pi}{6}$是$f\left(\dfrac{\pi}{2} - x\right) $的对称轴，$k = 3$对应$C.$
\end{jieda}

\item 设 $n$ 是正整数，集合 $A=\left\{x \left\lvert\, x=\cos \frac{2 k \pi}{n}\right., k \in \mathrm{Z}\right\}$ ．当 $n=2024$ 时，集合 A 元素的个数为\dotfill(\kh{B})
\xx
{1012}
{1013}
{2023}
{2024}

\item 已知函数$f(x)=\cos(2x-\frac{\pi}{3})+2\sin(x-\frac{\pi}{4})\sin(x+\frac{\pi}{4})$. 
\begin{enumerate}[(1)]
\item 求函数$f(x)$的最小正周期和图像的对称轴方程; 
\item 求函数$f(x)$在区间$[-\frac{\pi}{12},\frac{\pi}{2}]$上的值域. 
\end{enumerate}
\begin{jieda}
$f(x)=\sin(2x-\frac{\pi}{6})$. 
\begin{enumerate}
\item $T=\pi$; 轴: $x=\frac{k\pi}{2}+\frac{\pi}{3}(k\zz)$. 
\item $[-\frac{\sqrt{3}}{2},1]$. 
\end{enumerate}
\end{jieda}


\item 设函数$f(x)=2\cos^2x+2\sqrt{3}\sin{x}\cos{x}+m(x\zr)$. 
\begin{enumerate}[(1)]
\item 化简函数$f(x)$的表达式, 并求函数的最小正周期;
\item 若$x\in[0,\frac{\pi}{2}]$, 是否存在实数$m$, 使函数$f(x)$的值域恰为$[\frac{1}{2},\frac{7}{2}]$? 若存在, 请求出$m$的取值; 若不存在, 请说明理由. 
\end{enumerate}
\begin{jieda}
\begin{enumerate}[(1)]
\item $f(x)=2\sin(2x+\frac{\pi}{6})+m+1=2\cos(2x-\frac{\pi}{3})+m+1$. 
\item $f(x)\in[m,m+3]$, 则$m=\frac{1}{2}\land m+3=\frac{7}{2}$, 故$m=\frac{1}{2}$.  
\end{enumerate}
\end{jieda}


\item 已知函数$f(x)=a\sin{x}-\frac{1}{2}\cos2x+a-\frac{3}{a}+\frac{1}{2}$, $a\zr$, $a\neq0$. 
\begin{enumerate}[(1)]
\item 若对任意$x\zr$都有$f(x)\leq 0$求$a$的取值范围; 
\item 若$a\geq 2$且存在$x\zr$使得$f(x)\leq 0$, 求$a$的取值范围. 
\end{enumerate}
\begin{jieda}
$f(x)=\sin^2x+a\sin{x}+a-\frac{3}{a}=(\sin{x}+\frac{a}{2})^2-\frac{a^2}{4}+a-\frac{3}{a}$. 
\begin{enumerate}[(1)]
\item 法一: 先缩小范围: 将$\sin{x}=0,\pm1$代入得$1-\frac{3}{a}\leq 0\land a-\frac{3}{a}\leq 0\land 1+2a-\frac{3}{a}\leq 0$, \\
故$a\in(0,1]$. \\
于是$-\frac{a}{2}\in[-\frac{1}{2},0)$, 故$f(x)_{\min}=-\frac{a^2}{4}+a-\frac{3}{a}\leq 0$. \\
解得$a\in(0,1]$. \\
法二: 实际上, 对于开口向上的二次函数而言, 在闭区间上的最大值必然是在区间的端点上取到, 那么本题只要$\sin{x}=\pm1$时成立即可. 

\item 只要$f(x)_{\min}\leq 0$即可. \\
因为$a\geq 2$, 故$-\frac{a}{2}\leq -1$. \\
此时$f(x)_{\min}=1-\frac{3}{a}(\sin{x}=-1)$, 故$1-\frac{3}{a}\leq 0\land a\geq 2$, 即$a\in[2,3]$. 
\end{enumerate}
\end{jieda}

\item 已知函数 $f(x)=\sin k x \sin ^k x+\cos k x \cos ^k x-\cos ^k 2 x$ ，（其中 $k \in \mathrm{~N}^*$ ， $x \in \mathrm{R}$ ）（1）当 $k=1$ 时，求函数 $f(x)$ 的严格递增区间；
（2）当 $k=1$ 时，求函数 $g(x)=\frac{f(x)}{a+f^2(x)}$ 在 $\left(0, \frac{\pi}{3}\right]$ 上的最大值（其中常数 $a>0$ ）； （3）若函数 $f(x)$ 为常值函数，求 $k$ 的值．
\begin{jieda}
\begin{enumerate}[(1)]
    \item 当 $k=1$ 时，$f(x)=\sin x \sin x+\cos x \cos x-\cos 2 x=1-\cos 2 x$由 $2 k \pi \leq 2 x \leq 2 k \pi+\pi, k \in Z$ ，得 $k \pi \leq x \leq k \pi+\frac{\pi}{2}, k \in Z$ ．故 $f(x)$ 的严格递增区间为 $\left[k \pi, k \pi+\frac{\pi}{2}\right], k \in Z$ ．

    \item 由（1）可知，当 $k=1$ 时，$f(x)=1-\cos 2 x=2 \sin ^2 x$ ，则 $g(x)=\frac{f(x)}{a+f^2(x)}=\frac{2 \sin ^2 x}{a+4 \sin ^4 x}$ ，令 $t=2 \sin ^2 x$ ，当 $x \in\left(0, \frac{\pi}{3}\right]$ 时，则 $2 x \in\left(0, \frac{2 \pi}{3}\right]$ ，所以 $\cos 2 x \in\left[-\frac{1}{2}, 1\right)$ ，则 $1-\cos 2 x \in\left(0, \frac{3}{2}\right]$ ，即 $t \in\left(0, \frac{3}{2}\right]$ ．\\
    于是 $g(x)=h(t)=\frac{1}{\frac{a}{t}+t}, t \in\left(0, \frac{3}{2}\right]$
（1）当 $a \in\left(0, \frac{9}{4}\right]$ 时，$h(t)=\frac{1}{\frac{a}{t}+t} \leq \frac{1}{2 \sqrt{\frac{a}{t} \cdot t}}=\frac{1}{2 \sqrt{a}}$ ，当且仅当 $t=\sqrt{a}$ 时，最大值为 $\frac{1}{2 \sqrt{a}}$ ；
（2）当 $a>\frac{9}{4}$ 时，$y=\frac{a}{t}+t$ 在 $\left(0, \frac{3}{2}\right]$ 上递减，则 $h(t)$ 在 $\left(0, \frac{3}{2}\right]$ 上是增函数，则当 $t=\frac{3}{2}$ 时，最大值为 $\frac{6}{4 a+9}$ ，
综上所述，$g(x)_{\text {max }}=\left\{\begin{array}{l}\frac{\sqrt{a}}{2 a}, 0<a \leq \frac{9}{4} \\ \frac{6}{4 a+9}, a>\frac{9}{4}\end{array}\right.$
    \item 由函数 $f(x)$ 为常值函数，令 $x=0$ ，则原式 $=0$ ，令 $x=\frac{\pi}{2}$ ，则原式 $=\sin \frac{k \pi}{2}-(-1)^k=0 \Rightarrow k=4 n-1$（ $n$ 为正整数）；令 $x=\frac{\pi}{k}$ ，则原式 $=-\cos ^k \frac{\pi}{k}-\cos ^k \frac{2 \pi}{k}=0$ ，即 $\cos ^k \frac{\pi}{k}=-\cos ^k \frac{2 \pi}{k}$ ，因为 $k=4 n-1$（ $n$ 为正整数），即 $k$ 为正奇数，所以 $\cos \frac{\pi}{k}=-\cos \frac{2 \pi}{k}$ ，即 $\cos \frac{\pi}{k}+\cos \frac{2 \pi}{k}=0$ ，则 $2 \cos ^2 \frac{\pi}{k}+\cos \frac{\pi}{k}-1=0$ ，解得 $\cos \frac{\pi}{k}=-1$ 或 $\cos \frac{\pi}{k}=\frac{1}{2}$ ，

又因为 $k=4 n-1$（ $n$ 为正整数），所以 $k=3$ ．
当 $k=3$ 时，原式为
$$
\begin{aligned}
& \sin 3 x \sin ^3 x+\cos 3 x \cos ^3 x-\cos ^3 2 x=\sin 3 x \sin x \sin ^2 x+\cos 3 x \cos x \cos ^2 x-\cos ^3 2 x \\
& =\sin 3 x \sin x\left(1-\cos ^2 x\right)+\cos 3 x \cos x\left(1-\sin ^2 x\right)-\cos ^3 2 x \\
& =\sin 3 x \sin x+\cos 3 x \cos x-\sin 3 x \sin x \cos ^2 x-\cos 3 x \cos x \sin ^2 x-\cos ^3 2 x \\
& =\cos 2 x-\sin x \cos x(\sin 3 x \cos x+\cos 3 x \sin x)-\cos ^3 2 x \\
& =\cos 2 x-\sin x \cos x \sin 4 x-\cos ^3 2 x=\cos 2 x-\sin ^2 2 x \cos 2 x-\cos ^3 2 x \\
& =\cos 2 x\left(1-\sin ^2 2 x\right)-\cos ^3 2 x=\cos ^3 2 x-\cos ^3 2 x=0 .
\end{aligned}
$$
所以当 $k=3$ 时，函数 $f(x)$ 为常值函数．
\end{enumerate}   
\end{jieda}

\end{enumerate}


C组：

\begin{enumerate}

\item 设函数$f(x)=3\sin{x}+2\cos{x}+1$, 若实数$a,b,c$使得$af(x)+bf(x-c)=1$对任意实数$x$恒成立, 则$\dfrac{b\cos{c}}{a}$=\blkda{$-1$}. 
\begin{jieda}
法一: 全部展开. 
\begin{align*}
E & = a(3\sin{x}+2\cos{x}+1)+b(3\sin(x-c)+2\cos(x-c)+1)\\
& =(3a+3b\cos{c}+2b\sin{c})\sin{x}+(2a-3b\sin{c}+2b\cos{c})\cos{x}+(a+b)
\end{align*}
对于$m\sin{x}+n\cos{x}(x\zr)$而言, 当且仅当$m=n$时, 才会为常数, 且为零. \\
于是$\begin{cases}
3a+3b\cos{c}+2b\sin{c}=0\\
2a-3b\sin{c}+2b\cos{c}=0\\
a+b=1
\end{cases}$, 
即$\begin{cases}
\frac{a}{-b\cos{c}}=1+\frac{2}{3}\tan{c}\\
\frac{a}{-b\cos{c}}=1-\frac{3}{2}\tan{c}
\end{cases}$, \\
于是$\tan{c}=0$, 故$\frac{b\cos{c}}{a}=-1$. \\
法二: $f(x)=\sqrt{13}\sin(x+\varphi)+1$, \\
则$af(x)+bf(x-c)=(\sqrt{13}a+\sqrt{13}b\cos{c})\sin(x+\varphi)-\sqrt{13}b\sin{c}\cos(x+\varphi)+a+b$. \\
于是$\begin{cases}
\sqrt{13}(a+b\cos{c})=0\\
\sqrt{13}b\sin{c}=0\\
a+b=1
\end{cases}$. \\
由第二式, 若$b=0$, 则$a=0$(第一式). 与$a+b=1$矛盾. 故$b\neq 0$. 于是$\sin{c}=0$. \\
代入第一式, 结合第三式, 只能$a=b=\frac{1}{2}$, 且$\cos{c}=-1$. \\
法三: 令$x=0,x=c,x=\frac{\pi}{2}$之类的也许可以. 
\end{jieda}


\item 已知当$x\in[0,1]$时, 不等式$x^2\cos\theta-x(1-x)+(1-x)^2\sin\theta>0$恒成立, 求$\theta$的取值范围. 
\begin{jieda}
从形式上来看的话, 基本上就是$x^2\cos\theta+(1-x)^2\sin\theta>x(1-x)$, 基本不等式差不多了. \\
于是$\begin{cases}
\cos\theta\geq 0\\
\sin\theta\geq 0\\
2\sqrt{\cos\theta\sin\theta}>1\\
\end{cases}$, 
解得$\theta\in\left(2k\pi+\dfrac{\pi}{12},2k\pi+\dfrac{5\pi}{12}\right)(k\zz)$. 
\end{jieda}

\end{enumerate}



\newpage


\section{函数$y=A\sin(\omega{x}+\varphi)+B$的图像与性质(1)}
\zsd 函数$y=A\sin(\omega{x}+\varphi)+B$的图像与变换 \\

A组：
\begin{enumerate}

\item 函数 $y=-2 \sin \left(\frac{\pi}{4}-\frac{x}{2}\right)$ 的振幅为\blkda[1]{2}


\item 函数 $y=2 \sin \left[2\left(x+\frac{1}{6}\right)\right]$ 的初始相位是\blkda[1]{$\dfrac{1}{3}$}


\item 将函数 $y=\sin x$ 图像上各点的横坐标缩小为原来的 $\frac{1}{4}$ ，纵坐标不变，再将所得图像向左平移 $\frac{\pi}{24}$ 个单位长度得到函数 $y=f(x)$ 的图像，则 $f(x)$ 的解析式为\blkda[2]{$\sin{(4x+\dfrac{\pi}{6})}$}．

\item 函数 $f(x)=\sin ^2 x+\sin 2 x$ 的值域为\blkda[4]{$\left[\dfrac{1-\sqrt{5}}{2},\dfrac{1+\sqrt{5}}{2}\right]$} ．


\item 函数 $f(x)=\sin ^2 x+2 \sin x \cos x+3 \cos ^2 x$ 的周期为\blkda[3]{$k\pi, k \in \mathbf{Z} \mbox{且} k \neq 0$}．


\item 为了得到函数 $y=\sin (4 x+2)$ 的图象，可以将函数 $y=\sin (4 x-6)$ 的图象\dotfill(\kh{D})
\xx
{向右平移 8 个单位长度}
{向左平移 8 个单位长度}
{向右平移 2 个单位长度}
{向左平移 2 个单位长度}


\item 函数 $y=\sin (\omega x+\varphi)(x \in \mathbf{R}, \omega>0,0 \leq \varphi<2 \pi)$ 的部分图象如图，则\dotfill(\kh{C})
\xx
{$\omega=\frac{\pi}{2}, \varphi=\frac{\pi}{4}$}
{$\quad \omega=\frac{\pi}{3}, \varphi=\frac{\pi}{6}$}
{$\quad \omega=\frac{\pi}{4}, \varphi=\frac{\pi}{4}$}
{$\omega=\frac{\pi}{4}, \varphi=\frac{5 \pi}{4}$}

\end{enumerate}

\begin{minipage}{\textwidth}
\centering
\includegraphics[scale=0.9]{./bodyxb/src/main42xbb2ch06ruijiandesanjiaobi15t7.jpg}
\end{minipage}



B组：
\begin{enumerate}

\item 为了得到函数 $y=\sin \left(2 x+\frac{\pi}{3}\right)$ 的图象，可以将函数 $y=\cos \left(2 x-\frac{2 \pi}{3}\right)$ 的图象\dotfill(\kh{B})
\xx
{向左平移 $\frac{\pi}{2}$ 个单位}
{向左平移 $\frac{\pi}{4}$ 个单位}
{向右平移 $\frac{\pi}{2}$ 个单位}
{向右平移 $\frac{\pi}{4}$ 个单位}


\item 将 $y=f(x)$ 的图像向左平移 $\frac{\pi}{4}$ 个单位，再向上平移 1 个单位之后，可得 $y=\sin 2 x$ 的图像，则 $f\left(\frac{\pi}{12}\right)$=\blkda[2]{$-\dfrac{\sqrt{3}}{2}-1$}.


\item 把函数 $y=\cos \left(3 x+\frac{\pi}{4}\right)$ 的图像适当变动就可以得到 $y=\sin (-3 x)$ 图像，这种变动可以是\dotfill(\kh{D})
\xx
{向右平移 $\frac{\pi}{4}$}
{向左平移 $\frac{\pi}{4}$}
{向右平移 $\frac{\pi}{12}$}
{向左平移 $\frac{\pi}{12}$}


\item 已知函数$f(x)=2\cos^2\omega{x}+2\sin\omega{x}\cos\omega{x}+1(x\zr,\omega>0)$的最小值正周期是$\dfrac{\pi}{2}$, 
则函数$f(x)$取最大值时的$x$的集合为\blkda[4.5]{$\{x\mid x=\frac{k\pi}{2}+\frac{\pi}{16},\zz\}$}. 
\begin{jieda}
$f(x)=\sqrt{2}\sin(2\omega{x}+\frac{\pi}{4})+2$. \\
于是$\frac{\pi}{2}=\frac{2\pi}{2\omega}$, 故$\omega=2$. \\
故$f(x)=\sqrt{2}\sin(4x+\frac{\pi}{4})+2$, 于是$4x+\frac{\pi}{4}=2k\pi+\frac{\pi}{2}$. 
\end{jieda}

\item 已知函数$f(x)=\sqrt{3}\sin\omega{x}+\cos\omega{x}(\omega>0)$, $y=f(x)$的图像与直线$y=2$的两个相邻交点的距离等于$\pi$, 
则$f(x)$的单调递增区间是\blkda[4.5]{$[k\pi-\frac{\pi}{3},k\pi+\frac{\pi}{6}](k\zz)$}. 
\begin{jieda}
$f(x)=2\sin(\omega{x}+\frac{\pi}{6})$, 则最大值为2. \\
于是$T=\pi$, 易得$\omega=2$, 则$f(x)=2\sin(2x+\frac{pi}{6})$. \\
解$2x+\frac{\pi}{6}\in[2k\pi-\frac{\pi}{2},2k\pi+\frac{\pi}{2}]\iff x\in[k\pi-\frac{\pi}{3},k\pi+\frac{\pi}{6}](k\zz)$. 
\end{jieda}

\item 函数$y=A\sin(\omega{x}+\varphi)(A>0,\omega>0,\abs{\varphi}<\frac{\pi}{2})$在$x\in(0,\frac{2\pi}{3})$内只取到一个最大值和一个最小值,  
且当$x=\frac{\pi}{12}$时, 函数的最大值为3; 当$x=\frac{7\pi}{12}$时, 函数的最小值为$-3$, 试求此函数的解析式为\blkda[3]{$y=3\sin(2x+\frac{\pi}{3})$}. 
\begin{jieda}
因为$\frac{\pi}{12},\frac{7\pi}{12}\in(0,\frac{2\pi}{3})$, 结合条件可得$\frac{7\pi}{12}-\frac{\pi}{12}=\frac{T}{2}$, 故$\omega=2$. \\
此时$2\cdot\frac{\pi}{12}+\varphi\in(-\frac{\pi}{3},\frac{2\pi}{3})$, 故$\sin(2\cdot\frac{\pi}{12}+\varphi)\in(-\frac{\sqrt{3}}{2},1]$, \\
故$\sin(2\cdot\frac{\pi}{12}+\varphi)=1$, 且$A=3$. \\
\ps 上述讨论是必要的, 否则可能出现$A=-3$, $2\cdot\frac{\pi}{12}+\varphi=\frac{3\pi}{2}$之类的情形. \\
解得$\varphi=\frac{\pi}{3},A=3$, 即$y=3\sin(2x+\frac{\pi}{3})$. 
\end{jieda}

\item 在平面直角坐标系$xOy$中, 求函数$f(x)=a\sin{ax}+\cos{ax}(a>0)$在一个最小正周期长的区间上的图像与函数$g(x)=\sqrt{a^2+1}$的图像所围成的封闭图形的面积为\blkda{$\frac{2\pi\sqrt{a^2+1}}{a}$}. 
\begin{jieda}
作图易得. 
\end{jieda}

\item 已知函数$y=2\sin(\omega{x}+\varphi)(\varphi>0)$在区间$[0,2\pi]$的图像如下: 

那么$\omega$为\dotfill(\kh{B})
\xx{1}{2}{$\frac{1}{2}$}{$\frac{1}{3}$}
\begin{jieda}
$T=\pi$. 
\end{jieda}
\begin{minipage}{\textwidth}
\centering
\includegraphics[scale=0.9]{./bodyxb/src/main42xbb2ch06ruijiandesanjiaobi15t8.jpg}
\end{minipage}



\item 若函数$f(x)=3\sin(\omega{x}+\varphi)$对任意实数$x$, 都有$f(\frac{\pi}{4}+x)=f(\frac{\pi}{4}-x)$, 则$f(\frac{\pi}{4})$等于\dotfill(\kh{D})
\xx{0}{3}{$-3$}{3或$-3$}
\begin{jieda}
$f(\frac{\pi}{4})=\min$或$\max$. 
\end{jieda}


\item 函数 $y=A \sin (\omega x+\varphi)\left(\omega>0,|\varphi|<\frac{\pi}{2}, x \in R\right)$ 的部分图象如图所示，则函数表达式为\dotfill(\kh{A})
\xx
{$y=-4 \sin \left(\frac{\pi}{8} x+\frac{\pi}{4}\right)$}
{$\quad y=4 \sin \left(\frac{\pi}{8} x-\frac{\pi}{4}\right)$}
{$y=-4 \sin \left(\frac{\pi}{8} x-\frac{\pi}{4}\right)$}
{$\quad y=4 \sin \left(\frac{\pi}{8} x+\frac{\pi}{4}\right)$}
\begin{jieda}
$T=\pi$. 
\end{jieda}
\begin{minipage}{\textwidth}
\centering
\includegraphics[scale=0.9]{./bodyxb/src/main42xbb2ch06ruijiandesanjiaobi15t10.jpg}
\end{minipage}





\item 已知函数 $g(x)=2 \sin 2 x$ ，函数 $f(x)=A \sin (\omega x+\varphi)$（其中 $A>0, \omega>0$ ， $\left.|\varphi|<\frac{\pi}{2}\right)$ 的部分图象如图所示，要得到 $f(x)$ 的图象，只需将函数 $g(x)$ 的图象\dotfill(\kh{A})
\xx
{向左平移 $\frac{\pi}{6}$ 个单位}
{向右平移 $\frac{\pi}{6}$ 个单位}
{向右平移 $\frac{\pi}{3}$ 个单位}
{向左平移 $\frac{\pi}{3}$ 个单位}
\begin{minipage}{\textwidth}
\centering
\includegraphics[scale=0.9]{./bodyxb/src/main42xbb2ch06ruijiandesanjiaobi15t11.jpg}
\end{minipage}





\item 已知函数$y=2\sin(2x+\frac{\pi}{3})$. 
\begin{enumerate}[(1)]
\item 求它的振幅、周期和初相; 
\item 用五点法作出它的图像. 
\end{enumerate}
\begin{jieda}
\begin{enumerate}[(1)]
\item $A=2$, $T=\pi$, $\varphi=\frac{\pi}{3}$. % 有人振幅4, 初相为$\sqrt{3}$. 
\item $2x+\frac{\pi}{3}=0,\frac{\pi}{2},\pi,\frac{3\pi}{2},2\pi$. (列表, 描点, 连线)\\
\begin{tabular}{|c|c|c|c|c|c|}\hline
$x$ & $-\frac{\pi}{6}$ & $\frac{\pi}{12}$ & $\frac{\pi}{3}$ & $\frac{7\pi}{12}$ & $\frac{5\pi}{6}$\\\hline
$2x+\frac{\pi}{3}$ & 0 & $\frac{\pi}{2}$ & $\pi$ & $\frac{3\pi}{2}$ & $2\pi$\\\hline
$y$ & 0 & 2 & 0 & $-2$ & 0\\\hline
\end{tabular}
\end{enumerate}
\end{jieda}


\item 已知函数$f(x)=\cos(2x-\frac{\pi}{3})+2\sin(x-\frac{\pi}{4})\sin(x+\frac{\pi}{4})$. 
\begin{enumerate}[(1)]
\item 求函数$f(x)$的图像的对称轴方程; 
\item 求函数$f(x)$在区间$[-\frac{\pi}{12},\frac{\pi}{2}]$上的值域. 
\end{enumerate}
\begin{jieda}
$f(x)=\sin(2x-\frac{\pi}{6})$
\begin{enumerate}[(1)]
\item $x=\frac{k\pi}{2}+\frac{\pi}{3}(k\zz)$
\item $[-\frac{\sqrt{3}}{2},1]$
\end{enumerate}
\end{jieda}

\item 已知函数$f(x)=\sin^2\omega{x}+\sqrt{3}\sin\omega{x}\sin(\omega{x}+\frac{\pi}{2})(\omega>0)$的最小正周期为$\pi$. 
\begin{enumerate}[(1)]
\item 求$\omega$的值; 
\item 求函数$f(x)$在区间$[0,\frac{2\pi}{3}]$上的取值范围. 
\end{enumerate}
\begin{jieda}
$f(x)=\sin(2\omega{x}-\frac{\pi}{6})+\frac{1}{2}$. 
\begin{enumerate}[(1)]
\item $\omega=1$. 
\item $[0,\frac{3}{2}]$. 
\end{enumerate}
\end{jieda}

\item 已知函数$y=A\sin(\omega{x}+\varphi)$在一个周期内的图像的最高点的坐标为$(\frac{5\pi}{24},2)$, 最低点的坐标为$(\tfrac{11\pi}{24},-2)$. 
\begin{enumerate}[(1)]
\item 求此函数的解析式;
\item 求此函数的单调增区间.
\end{enumerate}
\begin{jieda}
\begin{enumerate}[(1)]
\item $f(x)=2\sin(4x-\frac{\pi}{3})$. 
\item $[\frac{k\pi}{2}-\frac{\pi}{24},\frac{k\pi}{2}+\frac{5\pi}{24}](k\zz)$. 
\end{enumerate}
\end{jieda}


\item 已知函数 $f(x)=2 \sin \left(2 x+\frac{\pi}{3}\right)$ ：
\begin{enumerate}[(1)]
    \item 若 $x \in\left[0, \frac{\pi}{4}\right]$ ，求 $y=f(x)$ 的最大值和最小值，并写出相应的 $x$ 值；
    \item 将函数 $y=f(x)$ 的图象向右平移 $\frac{\pi}{12}$ 个单位，再向上平移 1 个单位，得到函数 $y=g(x)$ 的图象，区间 $[a, b](a, b \in R$ 且 $a<b)$ 满足：$y=g(x)$ 在 $[a, b]$上至少含有 20 个零点，在所有满足上述条件的 $[a, b]$ 中，求 $b-a$ 的最小值．
\end{enumerate}
\begin{jieda}
    \begin{enumerate}[(1)]
        \item $2x+\dfrac{\pi}{3} \in \left[\dfrac{\pi}{3}, \dfrac{5\pi}{6}\right]$，故当$x=\dfrac{\pi}{12}$时，取$f(x)$最大值2；当$x=\dfrac{\pi}{4}$时，$f(x)$取最小值1
        \item $g(x) = f\left(x - \dfrac{\pi}{12}\right) + 1 = 2sin(2x+\dfrac{\pi}{6}) + 1$，$g(x)$的零点为$k\pi - \dfrac{\pi}{6} \mbox{或} k\pi -\dfrac{\pi}{2}, k \in \mathbf{Z}$，易知$g(x)$在一个周期内有两个零点，因此至少含有二十个零点时，区间长度$b-a$的最小值应为9个周期长度再加上两个零点间的最小距离，即$9\pi + \dfrac{\pi}{3} = \dfrac{28\pi}{3}$，当$a = k\pi -\dfrac{\pi}{2}, b = k\pi -\dfrac{\pi}{2} + \dfrac{28\pi}{3}，k \in \mathbf{Z}$时，可以囊括恰好20个零点.
    \end{enumerate}
\end{jieda}

\end{enumerate}

\section{函数$y=A\sin(\omega{x}+\varphi)+B$的图像与性质(2)}

\zsd 函数$y=A\sin(\omega{x}+\varphi)+d$的图像与变换 \\
A组：
\begin{enumerate}
\item 函数 $y=2 \sin \left[2\left(x+\frac{1}{6}\right)\right]$ 的初始相位是\blkda[1]{$\dfrac{1}{3}$}

\item 为了得到函数 $f(x)=\sin x(\cos x-\sqrt{3} \sin x)+\frac{\sqrt{3}}{2}$ 的图象，只需把函数 $f(x)=\sin 2 x$ 的图象上所有的点向\blkda{左}平移\blkda{$\frac{\pi}{6}$}个单位．

\item 将函数 $y=\sin x$ 图象上各点的横坐标缩小为原来的 $\frac{1}{4}$ ，纵坐标不变，再将所得图象向左平移 $\frac{\pi}{24}$ 个单位长度得到函数 $y=f(x)$ 的图象，则 $f(x)$ 在 $\left[0, \frac{\pi}{4}\right]$ 上的值域为\blkda{[$-\frac{1}{2},1$]}．

\item 函数 $f(x)=\cos ^2 x+\sqrt{3} \sin x \cos x$ 的值域为\blkda{[$-\frac{1}{2},\frac{3}{2}$]}.

\item 函数 $f(x)=\sqrt{3} \sin x+\sin \left(x+\frac{\pi}{2}\right)$ 的最大值是\blkda{$2$}．

\item 函数 $f(x)=\cos \left(2 x-\frac{\pi}{3}\right)+2 \sin \left(x-\frac{\pi}{4}\right) \sin \left(x+\frac{\pi}{4}\right)$ 最小正周期为\blkda{$\pi$}．
\end{enumerate}


B组：

\begin{enumerate}
    \item 将函数 $y=\sin 2 x$ 的图象向左平移 $\frac{\pi}{4}$ 个单位，再向上平移 1 个单位，所得图象的函数解析是\blkda[3]{$y=\cos 2x+1$}.

    \item 函数 $y=A \sin (\omega x+\varphi)\left(A>0, \omega>0,|\varphi|<\frac{\pi}{2}\right)$ 在 $x \in\left(0, \frac{2 \pi}{3}\right)$ 内只取到一个最大值和一个最小值，且当 $x=\frac{\pi}{12}$ 时，函数的最大值为 3 ，当 $x=\frac{7 \pi}{12}$ 时，函数的最小值为 $-3$ ，试求此函数的解析式为\blkda[3]{$y=3\sin(2x+\frac{\pi}{3})$}．

    \item 函数 $y=x+4+\sqrt{5-x^2}$ 的最大值是\blkda{$4+\sqrt{10}$}，最小值是\blkda{$4-\sqrt{5}$}．

    \item 简谐振动 $x_1=\sin \left(\omega t+\frac{\pi}{3}\right)$ 和 $x_2=\sin \left(\omega t-\frac{\pi}{6}\right)$ 叠加后得到的合振动是 $x=$ \blkda[3]{$\sqrt{2}\sin(\omega t+\frac{\pi}{2})$}．

    \item 已知函数 $f(x)=\sin \left(\omega x-\frac{\pi}{6}\right)$ ，其中 $\omega>0$ ，若 $f(x)$ 在区间 $\left(0, \frac{\pi}{3}\right)$ 上恰有 2 个零点，则 $\omega$ 的取值范围是\blkda{($\frac{7}{2},\frac{13}{2}$]}.

    \item 已知函数 $f(x)=\sin \left(\omega x+\frac{\pi}{4}\right)(x \in R, \omega>0)$ 的最小正周期为 $\pi$ ，为了得到函数 $g(x)=\cos \omega x$ 的图象，只要将 $y=f(x)$ 的图象\dotfill(\kh{A})
    \xx{向左平移 $\frac{\pi}{8}$ 个单位长度}{向右平移 $\frac{\pi}{8}$ 个单位长度}{向左平移 $\frac{\pi}{4}$ 个单位长度}{向右平移 $\frac{\pi}{4}$ 个单位长度}

    \item 为得到函数 $y=\cos \left(2 x+\frac{\pi}{3}\right)$ 的图像，只需将函数 $y=\sin 2 x$ 的图像\dotfill(\kh{A})
    \xx{向左平移 $\frac{5 \pi}{12}$ 个长度单位}{向右平移 $\frac{5 \pi}{12}$ 个长度单位}{向左平移 $\frac{5 \pi}{6}$ 个长度单位}{向右平移 $\frac{5 \pi}{6}$ 个长度单位}

    \item 将函数 $y=\sin x$ 的图象向左平移 $\varphi(0 \leq \varphi<2 \pi)$ 个单位后得到函数 $y=\sin \left(x-\frac{\pi}{6}\right)$ 的图象，则 $\varphi$ 为\blkda{$\frac{11}{6}\pi$}.

    \item 将函数 $f(x)=\sin (\omega x+\varphi)$ 的图像向左平移 $\frac{\pi}{2}$ 个单位．若所得图象与原图象重合，则 $\omega$ 的值不可能等于\dotfill(\kh{B})
    \xx{$4$}{$6$}{$8$}{$12$}

\item 已知函数$f(x)=A\sin(x+\varphi)(A>0, 0<\varphi<\pi)$, $x\zr$的最大值是1, 其图像经过点$M\left(\frac{\pi}{3},\frac{1}{2}\right)$. 
\begin{enumerate}
\item 求$f(x)$的解析式; 
\item 已知$\alpha,\beta\in\left(0,\frac{\pi}{2}\right)$, 且$f(\alpha)=\frac{3}{5}$, $f(\beta)=\frac{12}{13}$, 求$f(\alpha-\beta)$的值. 
\end{enumerate}
\begin{jieda}
\begin{enumerate}[(1)]
\item $f(x)=\sin(x+\frac{\pi}{2})$. 
\item $\frac{56}{65}$. 
\end{enumerate}
\end{jieda}

\item 已知函数$f(x)=A\sin(\omega{x}+\varphi), x\zr$(其中$A>0, \omega>0, 0<\phi<\frac{\pi}{2})$的图像与$x$轴的交点中, 
相邻两个交点之间的距离为$\frac{\pi}{2}$, 且图像上一个最低点为$M(\frac{2\pi}{3},-2)$. 
\begin{enumerate}[(1)]
\item 求$f(x)$的解析式; 
\item 当$x\in[\frac{\pi}{12},\frac{\pi}{2}]$, 求$f(x)$的值域.
\end{enumerate}
\begin{jieda}
\begin{enumerate}[(1)]
\item $f(x)=2\sin(2x+\frac{\pi}{6})$. 
\item $[-1,2]$. 
\end{enumerate}
\end{jieda}

\item 已知函数$f(x)=\frac{1}{2}\sin2x\sin\varphi+\cos^2x\cos\varphi-\frac{1}{2}\sin\left(\frac{\pi}{2}+\varphi\right)\left(0<\varphi<\pi\right)$, 其图像过点$(\frac{\pi}{6},\frac{1}{2})$. 
\begin{enumerate}[(1)]
\item 求$\varphi$的值; 
\item 将函数$y=f(x)$的图像上各点的横坐标缩短到原来的$\frac{1}{2}$, 纵坐标不变, 得到函数$y=g(x)$的图像, 求函数$g(x)$在$[0,\frac{\pi}{4}]$上的最大值和最小值. 
\end{enumerate}
\begin{jieda}
$f(x)=\frac{1}{2}\cos(2x-\varphi)$. 
\begin{enumerate}[(1)]
\item 由$f(\frac{\pi}{6})=\frac{1}{2}$, 得$\cos(\frac{\pi}{3}-\varphi)=1$. 又$0<\varphi<\pi$, 故$\varphi=\frac{\pi}{3}$. 
\item $g(x)=f(2x)=\frac{1}{2}\cos(4x-\frac{\pi}{3})$. \\
$\min=g(\frac{\pi}{4})=-\frac{1}{4}$, $\max=g(\frac{\pi}{12})=\frac{1}{2}$. 
\end{enumerate}
\end{jieda}

\item 设函数 $y=f(x)$ 的表达式为 $f(x)=\sin \left(2 x+\frac{\pi}{6}\right), x \in \mathbf{R}$ ．
\begin{enumerate}
    \item 求函数 $y=f(x)$ 在 $x \in[0, \pi]$ 上的单调递增区间；
    \item 将函数 $y=f(x)$ 的图像向左平移 $\frac{\pi}{6}$ 个单位，得到函数 $g(x)$ 的图像，函数 $h(x)=f(x)+g(x)$ ，当 $x_1, x_2 \in\left(\frac{\pi}{4}, \frac{5 \pi}{6}\right)$ ，且 $x_1 \neq x_2$ 时，有 $h\left(x_1\right)=h\left(x_2\right)$ ，求 $h\left(x_1+x_2\right)$ 的值．
\end{enumerate}  
\begin{jieda}
\begin{enumerate}
    \item 单增区间为 $\left[0, \frac{\pi}{6}\right]$ 和 $\left[\frac{2 \pi}{3}, \pi\right]$
    \item $g(x)=\sin \left(2\left(x+\frac{\pi}{6}\right)+\frac{\pi}{6}\right)=\sin \left(2 x+\frac{\pi}{2}\right)=\cos 2 x$, \\ $h(x)=\sin \left(2 x+\frac{\pi}{6}\right)+\cos 2 x=\frac{\sqrt{3}}{2} \sin 2 x+\frac{3}{2} \cos 2 x=\sqrt{3} \sin \left(2 x+\frac{\pi}{3}\right)$, \\
    $h\left(x_1+x_2\right)=\sqrt{3} \sin \left(\frac{7}{3} \pi+\frac{\pi}{3}\right)=\frac{3}{2}$
\end{enumerate}
\end{jieda}

\end{enumerate}

C组：

\begin{enumerate}
    \item 已知函数 $f(x)=\frac{\sqrt{3}}{2} \sin 2 x-\frac{1}{2}\left(\cos ^2 x-\sin ^2 x\right)-1(x \in \mathbf{R})$ ，将函数 $f(x)$ 的图像向左平移 $\frac{\pi}{6}$ 个单位后得到函数 $g(x)$ 的图像，设 $\triangle A B C$ 三个角 $A, B, C$ 的对边分别为 $a, b, c$ ．
    \begin{enumerate}
        \item 若 $c=\sqrt{7}, f(C)=0, \sin B=3 \sin A$ ，求 $a, b$ 的值；
        \item 若 $g(B)=0, m=(\cos A, \cos B), n=(1, \sin A-\cos A \tan B)$ ，求 $m \cdot n$ 的取值范围．
    \end{enumerate}
    \begin{jieda}
        $f(x)=\frac{\sqrt{3}}{2} \sin 2 x-\frac{1}{2} \cos 2 x-1=\sin \left(2 x-\frac{\pi}{6}\right)-1$,
        $g(x)=\sin \left(2\left(x+\frac{\pi}{6}\right)-\frac{\pi}{6}\right)-1=\sin \left(2 x+\frac{\pi}{6}\right)-1$
        \begin{enumerate}
            \item 
            由 $f(c)=0$ 得 $C=\frac{\pi}{3}$ 
            由 \ $\sin B=3 \sin A$ 得 $b=3 a$ .
            于是 $\cos C=\frac{a^2+b^2-c^2}{2 a b}=\frac{10 a^2-7}{6 a^2}=\frac{1}{2} $
            解得 $a=1$ . 那么 $b=3$
            \item $g(B)=0$ ．得 $B=\frac{\pi}{6}$ \\
            $ \vec{m} \cdot \vec{n} =\cos A+\cos B(\sin A-\cos A \tan B) =\cos A+\sin (A-B) =\sin \left(A+\frac{\pi}{6}\right) $ \\
            由$A \in\left(0, \frac{5}{6} \pi\right)$, 知  $\vec{m} \cdot \vec{n} \in(0,1]$.

        \end{enumerate}
    \end{jieda}

\item 若 $\tan x_1 \tan x_2 \cdots \tan x_{2012}=1$ ，则 $\sin x_1 \sin x_2 \cdots \sin x_{2012}$ 的最大值是 \blkda{$\frac{1}{2^{1006}}$}．

\item 已知函数 $f(x)=2\left(\sin ^4 x+\cos ^4 x\right)+m(\sin x+\cos x)^4$ 在 $x \in\left[0, \frac{\pi}{2}\right]$ 上有最大值 5 ，求实数 $m$ 的值．
\begin{jieda}
$f(x)=2\left(\sin ^2 x+\cos ^2 x\right)^2-4 \sin ^2 x \cos ^2 x+m(\sin x+\cos x)^1$
$=2-(2 \sin x \cos x)^2+m(\sin x+\cos x)^4$ \\
令 $t=\sin x+\cos x=\sqrt{2} \sin \left(x+\frac{\pi}{4}\right) \in[1, \sqrt{2}]$ ，则 $2 \sin x \cos x=t^2-1$ ，从而
$$f(x)=2-\left(t^2-1\right)^2+m t^4=(m-1) t^4+2 t^2+1 .$$
令 $u=t^2 \in[1,2]$ ，由题意知 $g(u)=(m-1) u^2+2 u+1$ 在区间 $[1,2]$ 上有最大值 $5$ ．\\
当 $m-1=0$ 时，$g(u)=2 u+1$ 在 $u=2$ 时有最大值 5 ，故 $m=1$ 符合条件；\\
当 $m-1>0$ 时，$g(u)_{\text {max }} \geqslant g(2)>2 \times 2+1=5$ ，矛盾！\\
当 $m-1<0$ 时，$g(u)<2 u+1 \leqslant 5$ ，矛盾！\\
综上所述，所求的实数 $m=1$ ．
\end{jieda}
\end{enumerate}



\section{函数$y=A\sin(\omega{x}+\varphi)+B$的图像与性质(3)}
\zsd 函数$y=A\sin(\omega{x}+\varphi)+B$的图像与变换 \\

A组：
\begin{enumerate}
    \item 函数$f(x) = sin\left(\dfrac{\pi x}{4}-\dfrac{\pi}{6}\right) - 2cos^2\left(\dfrac{\pi x}{8}\right) + 1$的最小正周期为\blkda{8}
    \begin{jieda}
        $f(x) = sin\left(\dfrac{\pi x}{4}-\dfrac{\pi}{6}\right) + cos\left(\dfrac{\pi x}{4}\right)$
    \end{jieda}
    \item 函数$y = Acos(\omega x + \varphi)\left(A < 0, \omega > 0, |\varphi| < \dfrac{\pi}{2}\right)$在一个周期内的图像如下图所示，则函数解析式为\blkda[4]{$y = -2cos\left(\dfrac{x}{2} + \dfrac{\pi}{4} \right)$}
    \begin{figure}[h]
        \centering
        \includegraphics[width=0.25\linewidth]{bodyxb//src/Asin(omega x+phi)+B(3)_A2.png}
    \end{figure}
    \begin{jieda}
        首先根据最大值判断出来$A = -2$，根据平衡位置和对称轴之间的距离可知$T = 4\pi$，因此$\omega = \dfrac{1}{2}$，因此$y = -2cos\left(\dfrac{x}{2}+\varphi\right)$，代入波峰的坐标，得到$\varphi = \dfrac{\pi}{4}$
    \end{jieda}
    \item 如图，已知函数$y = Acos(\omega x + \varphi) (A > 0, \omega > 0, 0 < \varphi < 2\pi)$的图像和$y$轴交点为$(0, 1)$，并已知其在$y$轴右侧的第一个最高点和第一个最低点的坐标分别为$(x_0, 2)$和$(x_0 + 2\pi, -2)$，则函数解析式为\blkda[4]{$y = 2cos\left(\dfrac{x}{2} + \dfrac{5\pi}{3}\right)$}
    \begin{figure}[h]
        \centering
        \includegraphics[width=0.25\linewidth]{bodyxb//src/Asin(omega x + phi)+B(3)_A3.png}
    \end{figure}
    \begin{jieda}
        首先可以判断出$A = 2$，然后根据两个对称轴之间的坐标差可知$T = 4\pi$，故$\omega = \dfrac{1}{2}$，故$y = 2cos\left(\dfrac{x}{2} + \varphi \right)$，代入$(0, 1)$可得$\varphi = \dfrac{5\pi}{3}$或$\dfrac{\pi}{3}$（由单调性舍）
    \end{jieda}
    \item 设函数$f(x) = 2sin(\omega x + \varphi) \left(\omega > 0, |\varphi| < \dfrac{\pi}{2}\right)$，若$f(x)$的图象经过点$(0, 1)$，且$f(x)$在$[0, \pi]$上恰有两个零点，则实数$\omega$的取值范围是\blkda{$\left[\dfrac{11}{6}, \dfrac{17}{6}\right)$}
    \begin{jieda}
        首先带入$(0, 1)$可知$\varphi = \dfrac{\pi}{6}$，故$f(x) = 2sin\left(\omega x + \dfrac{\pi}{6}\right)$ \\
        令$t = \omega x + \dfrac{\pi}{6}$，故$f(x) = g(t) = 2sint$在$\left[\dfrac{\pi}{6}, \dfrac{\pi}{6} + \omega \pi\right]$上恰有两个零点，因为$g(t)$的零点为$k\pi, k \in \mathbf{Z}$，因此$k\pi \in \left[\dfrac{\pi}{6}, \dfrac{\pi}{6} + \omega \pi\right], k \in \mathbf{Z}$有两解，即$\left \{x | x = k - \dfrac{1}{6}, k \in \mathbf{Z}\right \}$与$[0, \omega]$有两个公共元素.\\
        结合数轴可知：两个公共元素应为$\dfrac{5}{6}, \dfrac{11}{6}$，$\omega \in \left[\dfrac{11}{6}, \dfrac{17}{6}\right)$
    \end{jieda}
\end{enumerate}
B组：
\begin{enumerate}
\item 把函数$y=\sqrt{3}\cos{x}-\sin{x}$的图像向左平移$m(m>0)$个单位, 所得的图像关于$y$轴对称, 则$m$的最小值为\blkda{$\dfrac{5\pi}{6}$}. 
\begin{jieda}
$f(x)=2sin(x+\dfrac{2\pi}{3})$，直线$x = m$是$f(x)$的对称轴，$f(x)$的对称轴为$x+\dfrac{2\pi}{3} = k\pi + \dfrac{\pi}{2}, k \in \mathbf{Z}$即$x = k\pi - \dfrac{\pi}{6}, k \in \mathbf{Z}$，故$m$最小值为$\dfrac{5\pi}{6}$
\end{jieda}

\item 已知函数$y=A\sin(\omega{x}+\varphi)+B(A>0, \omega>0, \abs{\varphi}<\frac{\pi}{2})$在同一周期中的最高点坐标为$(2,2)$, 
最低点坐标为$(8,-4)$, 则$y$的表达式为\blkda[4]{$y=3\sin(\frac{\pi}{6}x+\frac{\pi}{6})-1$}. 

\item 若函数$y=f(x)$的图像上每一点的纵坐标保持不变, 横坐标伸长到原来的2倍, 然后将所得图像先向左平移$\frac{\pi}{2}$ 个单位, 
再向下平移1个单位, 得到的曲线与$y=\frac{1}{2}\cos{x}$的图像相同, 则$y=f(x)$的表达式为\blkda[4]{$f(x)=\frac{1}{2}\sin2x+1$}. 

\item 设$\omega>0$, 函数$y=\sin(\omega{x}+\frac{\pi}{3})+2$的图像向右平移$\frac{4\pi}{3}$个单位后与原图像重合, 则$\omega$的最小值是\blkda{$\frac{3}{2}$}.
\begin{jieda}
$\frac{4\pi}{3}=kT(k\zz, k \neq 0)$. 
\end{jieda}

\item 如图，已知 $A,B$ 分别是函数 $f(x) = \sqrt{3} sin\omega x(\omega > 0)$在 $y$
轴右侧图象上的第一个最高点和第一个最低点，且$\angle AOB = \dfrac{\pi}{2}$，则该函数的周期是\blkda[4]{$4k, k \in \mathbf{Z}$且$k \neq 0$}
\begin{figure}[h]
    \centering
    \includegraphics[width=0.25\linewidth]{bodyxb//src/Asin(omega x + phi)+B(3)_B5.png}
\end{figure}
\begin{jieda}
    函数最小正周期$T = \dfrac{2\pi}{\omega}$，$A\left(\dfrac{T}{4}, \sqrt{3}\right), B\left(\dfrac{3T}{4}, -\sqrt{3}\right)$，因为$\angle AOB = \dfrac{\pi}{2}$，所以$\dfrac{3T^2}{16} = 3$，故$T = 4$
\end{jieda}
\item 为了得到函数$y=\sin(3x+\frac{\pi}{6})$的图像, 只需把函数$y=\sin3x$的图像\dotfill(\kh{B})
\xx
{向左平移$\frac{\pi}{6}$}
{向左平移$\frac{\pi}{18}$}
{向右平移$\frac{\pi}{6}$}
{向右平移$\frac{\pi}{18}$}


\item 把函数$y=-3\cos(2x+\frac{\pi}{3})$的图像向右平移$m(m>0)$个单位, 设所得图像的解析式为$y=f(x)$, 则当$y=f(x)$是偶函数时, $m$的值可以是\dotfill(\kh{B})
\xx
{$\frac{\pi}{3}$}
{$\frac{\pi}{6}$}
{$\frac{\pi}{4}$}
{$\frac{\pi}{12}$}


\item 函数$f(x)=\sin^2x+\sqrt{3}\sin{x}\cos{x}$在区间$[\frac{\pi}{4},\frac{\pi}{2}]$上的最大值是\blkda{$\dfrac{3}{2}$}
\begin{jieda}
$f(x)=\sin(2x-\frac{\pi}{6})+\frac{1}{2}$. \\
$x\in[\frac{\pi}{4},\frac{\pi}{2}]$, 则$2x-\frac{\pi}{6}\in[\frac{\pi}{3},\frac{5\pi}{6}]$, 故$f(x)_{\max}=\frac{3}{2}$. 
\end{jieda}

\item 已知函数$y = f(x)$且$f(x) = \left\{\begin{array}{l}
     sinx, sinx \geq cosx  \\
     cosx, sinx < cosx
\end{array}\right.$，给出下列三个命题：
\begin{dingautolist}{172}
    \item 函数的值域为$[-1, 1]$
    \item 当且仅当$2k\pi + \pi < x < 2k\pi + \dfrac{3\pi}{2}, k \in \mathbf{Z}$时$f(x) < 0$
    \item 对任意$x \in \mathbf{R}$，$f^2(x) + f^2(\pi + x) = 1$恒成立
\end{dingautolist}
上述命题中正确的序号是\blkda{\ding{173}\ding{174}}
\begin{jieda}
    \begin{dingautolist}{172}
        \item 显然函数值取不到$-1$，故错误
        \item 在第三象限的终边的正弦值和余弦值都是负的，因此$f(x) < 0$，而对于其他象限的终边，正弦值和余弦值至少有一个为正，因此$f(x) > 0$，故正确
        \item 终边旋转半圈后，正弦和余弦的大小关系改变，结合诱导公式和平方关系可知正确
    \end{dingautolist}
\end{jieda}
\item 已知函数$f(x) = Asin(\omega x + \varphi) \left(A > 0, \omega > 0, |\varphi| < \dfrac{\pi}{2} \right)$的部分图像如图所示，与$x$轴的交于点$A\left(-\dfrac{\pi}{12}, 0\right)$，且平行四边形$EDCB$的面积为$4\pi$；若函数$f(x)$在区间$\left(\dfrac{2\pi}{3}, m\right)$上单调递增，则实数$m$的取值范围是\blkda{$\left(\dfrac{2\pi}{3}, \dfrac{7\pi}{6}\right]$}
\begin{figure}[h]
    \centering
    \includegraphics[width=0.25\linewidth]{bodyxb//src/Asin(omega x + phi) + B(3)_B10.png}
\end{figure}
\begin{jieda}
    根据最大最小值可知$A = 2$，根据平行四边面积可知$T = \pi$，即$\omega = 2$，根据$A\left(-\dfrac{\pi}{12}, 0\right)$可知$E\left(\dfrac{\pi}{6}, 2\right)$，因此$\varphi = \dfrac{\pi}{6}$，即$f(x) = 2sin(2x+\dfrac{\pi}{6})$，其递增区间为$\left[k\pi - \dfrac{\pi}{3}, k\pi + \dfrac{\pi}{6}\right], k \in \mathbf{Z}$，故$m \in \left(\dfrac{2\pi}{3}, \dfrac{7\pi}{6}\right]$ \\
    \ps 得到$D\left(\dfrac{2\pi}{3}, -2\right), B\left(\dfrac{7\pi}{6}, 2\right)$判断亦可.
\end{jieda}
\item 如图是函数$f(x) = Asin(\omega x + \varphi), x \in \mathbf{R}\left(A > 0, \omega > 0, 0 < \varphi < \dfrac{\pi}{2} \right)$的部分图像，$M, N$是它与$x$轴的两个不同交点，$D$是$M,N$之间的最高点且横坐标为$\dfrac{\pi}{4}$，点$F(0, 1)$是线段$M, N$的中点.
\begin{figure}[h]
    \centering
    \includegraphics[width=0.25\linewidth]{bodyxb//src/Asin(omega x + phi)+B(3)_B11.png}
\end{figure}
\begin{enumerate}
    \item 求$f(x)$的解析式
    \item 若$x \in \left[-\dfrac{\pi}{12}, \dfrac{5\pi}{12}\right]$时，函数$h(x) = f(x) - 2a$有两个零点，求实数$a$的取值范围.
\end{enumerate}
\begin{jieda}
    \begin{enumerate}
        \item 根据几何关系可以写出$D\left(\dfrac{\pi}{4}, 2\right), M\left(-\dfrac{\pi}{4}, 0\right), N\left(\dfrac{3\pi}{4}, 0\right)$，故$A = 2, T = 2\pi, \omega = 1, \varphi = \dfrac{\pi}{4}$，故$f(x) = 2sin\left(x+\dfrac{\pi}{4}\right)$
        \item $h(x)$有两个零点即$sin\left(x+\dfrac{\pi}{4}\right) = a$有两解，$x \in \left[-\dfrac{\pi}{12}, \dfrac{5\pi}{12}\right], x+\dfrac{\pi}{4} \in \left[\dfrac{\pi}{6}, \dfrac{2\pi}{3}\right]$，即$a = sint$在$t \in \left[\dfrac{\pi}{6}, \dfrac{2\pi}{3}\right]$有两解，因此$a \in \left[\dfrac{\sqrt{3}}{2}, 1\right)$
    \end{enumerate}
\end{jieda}
\item 已知函数$f(x)=\sin(\omega{x}+\varphi)$, 其中$\omega>0$, $\abs{\varphi}<\frac{\pi}{2}$. 
\begin{enumerate}[(1)]
\item 若$\cos\frac{\pi}{4}\cos\varphi-\sin\frac{3\pi}{4}\sin\varphi=0$, 求$\varphi$的值; 
\item 在(1)的条件下, 若函数$f(x)$的图像的相邻两条对称轴之间的距离等于$\frac{\pi}{3}$, 求函数$f(x)$的解析式; 
并求最小正实数$m$, 使得函数$f(x)$的图像象左平移$m$个单位所对应的函数是偶函数. 
\end{enumerate}
\begin{jieda}
\begin{enumerate}[(1)]
\item $\cos(\varphi+\frac{\pi}{4})=0$, $\varphi\in(-\frac{\pi}{4},\frac{\pi}{4})$, 故$\varphi=\frac{\pi}{4}$. 
\item $\frac{\pi}{3}=\frac{T}{2}$, 故$\omega=3$. \\
于是$f(x)=\sin(3x+\frac{\pi}{4})$. \\
$3m+\frac{\pi}{4}=k\pi+\frac{\pi}{2}$, 故$m=\frac{k\pi}{3}+\frac{\pi}{12}$, 故$m_{\min}=\frac{\pi}{12}$. 
\end{enumerate}
\end{jieda}
\item 已知函数$f(x)=\sqrt{3}\sin(\omega{x}+\varphi)-\cos(\omega{x}+\varphi)(0<\varphi<\pi, \omega >0)$为偶函数, 
且函数$y=f(x)$图像的两相邻对称轴间的距离为$\frac{\pi}{2}$. 
\begin{enumerate}[(1)]
\item 求$f(\frac{\pi}{8})$的值; 
\item 将函数$y=f(x)$的图像向右平移$\frac{\pi}{6}$个单位后, 再将得到的图像上各点的横坐标伸长到原来的4倍, 纵坐标不变, 得到函数$y=g(x)$的图像, 求$g(x)$的单调递减区间. 
\end{enumerate}
\begin{jieda}
$f(x)=2\sin(\omega{x}+\varphi-\frac{\pi}{6})$. 
\begin{enumerate}[(1)]
\item $\varphi-\frac{\pi}{6}=k\pi+\frac{\pi}{2}$, $\varphi\in(0,\pi)$, 故$\varphi=\frac{2\pi}{3}$. \\
$\frac{\pi}{2}=\frac{T}{2}$, 故$\omega=2$. \\
于是$f(x)=2\sin(2x+\frac{\pi}{2})=2\cos2x$, 故$f(\frac{\pi}{8})=\sqrt{2}$. 

\item $g(x)=f(\frac{1}{4}x-\frac{\pi}{6})=2\cos(\frac{x}{2}-\frac{\pi}{3})$. \\
易得单调递减区间为$[4k\pi+\frac{2\pi}{3},4k\pi+\frac{8\pi}{3}](k\zz)$. 
\end{enumerate}
\end{jieda}

\item 已知函数$f(t)=\sqrt{\dfrac{1-t}{1+t}}$, $g(x)=\cos{x}\cdot f(\sin{x})+\sin{x}\cdot f(\cos{x})$, $x\in(\pi,\frac{17\pi}{12}]$. 
\begin{enumerate}[(1)]
\item 将函数$g(x)$化简成$A\sin(\omega{x}+\varphi)+B(A>0,\omega>0,\varphi\in[0,2\pi))$的形式; 
\item 求函数$g(x)$的值域. 
\end{enumerate}
\begin{jieda}
\begin{enumerate}[(1)]
\item $g(x)=\sqrt{2}\sin(x+\frac{\pi}{4})-2$. 
\item $g(x)\in[-\sqrt{2}-2,-3)$. 
\end{enumerate}
\end{jieda}
\item 已知函数$f(x)=a\sin\omega{x}+b\cos\omega{x}$($a,b,\omega\zr$, 且$\omega>0$)的部分图像如图所示. \\
\begin{minipage}{\textwidth}
\centering
\includegraphics[totalheight=6\baselineskip]{./bodyxb/src/main42xbb2ch07sanjiaohanshu01.png}
\end{minipage}
\begin{enumerate}[(1)]
\item 求$a,b,\omega$的值; 
\item 若方程$3(f(x))^2-f(x)+m=0$在$x\in(-\frac{\pi}{3},\frac{2\pi}{3})$内有两个不同的解, 求实数$m$的取值范围. 
\end{enumerate}
\begin{jieda}
\begin{enumerate}[(1)]
\item $\frac{7\pi}{6}-\frac{2\pi}{3}=\frac{T}{4}$, 且$\omega>0$, 故$\omega=1$. \\
$f(\frac{7\pi}{6})=-\sqrt{a^2+b^2}\land a^2+b^2=1$, 故$a=\frac{1}{2}$, $b=\frac{\sqrt{3}}{2}$. \\
综上, $f(x)=\sin(x+\frac{\pi}{3})$. 

\item 注意到$x\in(-\frac{\pi}{3},\frac{2\pi}{3})$, 对于方程$t=f(x)$而言, 当$t\leq 0$或$t>1$时, 无解; 当$t\in(0,1)$时, 有两解; 当$t=1$时, 有1解. \\
对$3t^2-t+m=0$, 则$t_1+t_2=\frac{1}{3}$, 且$t=\frac{1}{6}\pm d(d\geq 0)$, 其中$d=\frac{\sqrt{1-12m}}{6}$. \\
作直线$y=\frac{1}{6}\pm d$. \\
当$d=0$时, 符合条件, 此时$m=\frac{1}{12}$; \\
当$d>0$时, $\frac{1}{6}-d\leq 0\land\frac{1}{6}+d\in(\frac{1}{6},1)$, 即$d\in[\frac{1}{6},\frac{5}{6})$. \\
解得$m\in(-2,0)$. \\
综上, $m\in(-2,0]\cup\{\frac{1}{12}\}$. 
\end{enumerate}
\end{jieda}
\end{enumerate}
\clearpage

\section{正切函数的图像与性质(1)}
\zsd 正切函数的图像与性质
A组：
\begin{enumerate}
\item 正切函数的最小正周期是\blkda{$\pi$}，对称中心为\blkda[3]{$\left(\dfrac{k\pi}{2}, 0\right), k \in \mathbf{Z}$}
\item 有人说: ``正切函数在整个定义域内是单调递增的函数.'' 这句话对吗? 为什么? \blkda{不对}. 
\begin{jieda}
反例: $\tan\frac{\pi}{3}>\tan(\pi+\frac{\pi}{6})$. 
\end{jieda}
\item 函数$tan(ax+b)(a \neq 0)$的最小正周期为\blkda{$\dfrac{\pi}{|a|}$}
\item 根据正切函数的图像，写出使下列不等式成立的$x$值的集合：\\
\begin{enumerate*}
    \item $1+tanx \geq 0$
    \item $tanx - \sqrt{3} \geq 0$
\end{enumerate*}
\begin{jieda}
    \begin{enumerate*}
    \item $\{x | -\dfrac{\pi}{4}+k\pi \leq x < \dfrac{\pi}{2}+k\pi, k \in \mathbf{Z}\}$
    \item $\{x | \dfrac{\pi}{3}+k\pi \leq x < \dfrac{\pi}{2}+k\pi, k \in \mathbf{Z}\}$
\end{enumerate*}
\end{jieda}
\end{enumerate}
B组：
\begin{enumerate}
    \item 函数$y=\tan2x$的单调递增区间为\blkda[4]{$(\frac{k\pi}{2}-\frac{\pi}{4},\frac{k\pi}{2}+\frac{\pi}{4})(k\zz)$}. 
    \begin{jieda}
    $2x\in(k\pi-\frac{\pi}{2},k\pi+\frac{\pi}{2})$
    \end{jieda}
    \item 函数$y=\tan{x}-\cot{x}$的最小正周期为\blkda{$\frac{\pi}{2}$}. 
    \begin{jieda}
    $y=\dfrac{\sin{x}}{\cos{x}}-\dfrac{\cos{x}}{\sin{x}}=\dfrac{\sin^2x-\cos^2x}{\sin{x}\cos{x}}=\dfrac{-\cos2x}{\frac{1}{2}\sin2x}=-2\cot2x$. \\
    \ps 以上是恒定变形(没有改变定义域). 
    \end{jieda}
    \item 函数$y=\frac{1}{1+\tan2x}$的定义域为\blkda[6]{$\{x\mid x\neq \frac{k\pi}{2}-\frac{\pi}{8}\text{且}x\neq\frac{k\pi}{2}+\frac{\pi}{4}, k\zz\}$}. 
    \begin{jieda}
    $\tan2x\neq -1\iff 2x\neq k\pi-\frac{\pi}{4}\land 2x\neq k\pi+\frac{\pi}{2}$. 
    \end{jieda}
    
    \item 已知$\frac{\pi}{3}<\theta<\frac{\pi}{2}$, 比较$\sin\theta,\cot\theta,\cos\theta$的大小: \blkda[4]{$\sin\theta>\cot\theta>\cos\theta$}. 
    \begin{jieda}
    六个三角比皆为正, 于是做变形就方便点. \\
    $\sin\theta>\cos\theta$显然的. \\
    而$\cot\theta=\dfrac{\cos\theta}{\sin\theta}$, 那么$\cot\theta>\cos\theta$稍微想一下即得. \\
    又$\cos\theta\in(0,\frac{1}{2})$, 则
    \begin{align*}
    \cot\theta-\sin\theta 
    & =\dfrac{\cos\theta-\sin^2\theta}{\sin\theta}\\
    & =\dfrac{\cos^2\theta+\cos\theta-1}{\sin\theta}\\
    & =\dfrac{(\cos\theta+\frac{1}{2})^2-\frac{5}{4}}{\sin\theta}\\
    & <0
    \end{align*}
    综上, $\sin\theta>\cot\theta>\cos\theta$. 
    \end{jieda}
    
    \item 函数$y=\abs{\tan{x}}$的最小正周期为\blkda{$\pi$}
    \item 函数$y=\sqrt{1-\tan^2x}$的定义域为\blkda[4]{$[k\pi-\frac{\pi}{4},k\pi+\frac{\pi}{4}](k\zz)$}
    \item 关于$x$的不等式$\tan2x\leq-1$的解集为\blkda[4]{$(\frac{k\pi}{2}-\frac{\pi}{4},\frac{k\pi}{2}-\frac{\pi}{8}](k\zz)$}
    \begin{jieda}
    $k\pi-\frac{\pi}{2}<2x\leq k\pi-\frac{\pi}{4}$. 
    \end{jieda}
    \item $f(x) = \left \{ \begin{array}{l}
         tan(\pi x), x \geq 0 \\
         f(x+1), x < 0
    \end{array} \right.$，则$f\left(-\dfrac{2}{3}\right) = $\blkda{$\sqrt{3}$}
    \begin{jieda}
        $f\left(-\dfrac{2}{3}\right) = f\left(\dfrac{1}{3}\right) = tan\left(\dfrac{\pi}{3}\right) = \sqrt{3}$
    \end{jieda}
    \item 已知函数$f(x) = tan\left(\omega x + \dfrac{\pi}{3}\right) (\omega > 0)$在区间$\left(-\dfrac{\pi}{3}, \dfrac{\pi}{3}\right)$上单调，则$\omega$的取值范围为\blkda{$\left(0, \dfrac{1}{2}\right]$}
    \begin{jieda}
        换元$t = \omega x + \dfrac{\pi}{3}$，故问题等价于$tant$在$\left((1-\omega) \cdot \dfrac{\pi}{3}, (1+\omega) \cdot \dfrac{\pi}{3}\right)$上单调递增，又因为$tant$在$\left(k\pi - \dfrac{\pi}{2}, k\pi+\dfrac{\pi}{2} \right), k \in \mathbf{Z}$上单调递增 \\
        因此$\left((1-\omega) \cdot \dfrac{\pi}{3}, (1+\omega) \cdot \dfrac{\pi}{3}\right) \subseteq \left(k\pi - \dfrac{\pi}{2}, k\pi+\dfrac{\pi}{2} \right), k \in \mathbf{Z}$，\\
        即$\left(1-\omega, 1+\omega\right) \subseteq \left(3k-\dfrac{3}{2}, 3k+\dfrac{3}{2}\right), k \in \mathbf{Z}$，易知$1 \in \left(3k-\dfrac{3}{2}, 3k+\dfrac{3}{2}\right), k \in \mathbf{Z}$，因此必有$k = 0$，故$\omega \in \left(0, \dfrac{1}{2}\right]$
    \end{jieda}
    \item 判断函数$f(x)=\lg\dfrac{\tan{x}+1}{\tan{x}-1}$的奇偶性. 
    \begin{jieda}
    $\dfrac{\tan{x}+1}{\tan{x}-1}>0\iff -1<\tan{x}<1$, 由于对称性, 故解集对称. \\
    于是$f(-x)+f(x)=\lg\dfrac{\tan(-x)+1}{\tan(-x)-1}+\lg\dfrac{\tan{x}+1}{\tan{x}-1}=\lg\dfrac{1-\tan^2x}{(-1)^2-\tan^2x}=0$, 故``奇''. 
    \end{jieda}

    \item 已知函数$y = f(x)$，其中$f(x) = Atan(\omega x + \varphi), (\omega > 0, |\varphi| < \frac{\pi}{2})$，$f(x)$的部分图像如下图.
    \begin{enumerate*}
        \item 求$A, \omega, \varphi$的值
        \item 求$y = f(x)$的单调增区间
    \end{enumerate*}
    \begin{figure}[h]
        \centering
        \includegraphics[width=0.25\linewidth]{bodyxb//src/zhengqietuxiang(1)B11.png}
    \end{figure}
    \begin{jieda}
        由图可知$T = \dfrac{\pi}{2}$，故$\omega = 2$，代入$\left(\dfrac{3\pi}{8}, 0\right)$可知$\varphi = \dfrac{\pi}{4}$，代入$(0, 1)$可知$A = 1$，故$f(x) = tan\left(2x+\dfrac{\pi}{4}\right)$，$f(x)$单调区间为$\left(\dfrac{k\pi}{2} + \dfrac{\pi}{8}, \dfrac{(k+1)\pi}{2} + \dfrac{\pi}{8}\right), k \in \mathbf{Z}$
    \end{jieda}

    \item 已知函数$f(x) = 4tanx \cdot sin\left(\dfrac{\pi}{2} - x\right) \cdot cos\left(x - \dfrac{\pi}{3}\right) - \sqrt{3}$
    \begin{enumerate*}
        \item 求$f(x)$定义域
        \item 化简并求$f(x)$最小正周期
        \item 讨论$f(x)$在区间$\left[-\dfrac{\pi}{4}, \dfrac{\pi}{4}\right]$上的单调性
    \end{enumerate*}
    \begin{jieda}
        $f(x)$定义域为$\{x| x \neq k\pi + \dfrac{\pi}{2}, k \in \mathbf{Z}\}$ \\
        $f(x) = 4sinx \cdot \left(\dfrac{cosx}{2} + \dfrac{\sqrt{3}sinx}{2}\right) - \sqrt{3} = sin2x - \sqrt{3}cos2x = 2sin\left(2x - \dfrac{\pi}{3} \right), x \neq k\pi + \dfrac{\pi}{2}$，最小正周期为$\pi$ \\
        $f(x)$在$\left[-\dfrac{\pi}{4}, -\dfrac{\pi}{12}\right]$严格单调递减，在$\left(-\dfrac{\pi}{12}, \dfrac{\pi}{4}\right]$严格单调递增
    \end{jieda}
    \item 求函数$y=\dfrac{\sin{x}}{\sin{x}+\cos{x}}(x\in[0,\tfrac{\pi}{2}))$的最值. 
    \begin{jieda}
    当$x=0$时, $y=0$; \\
    当$x\in(0,\frac{\pi}{2})$时, $y=\dfrac{1}{1+\cot{x}}\in(0,1)$. \\
    综上, $y\in[0,1)$, 即最小值为0, 无最大值. 
    \end{jieda}

    \item 求函数$y=\dfrac{\tan^2x-\tan{x}+1}{\tan^2x+\tan{x}+1}$的最大值与最小值, 并求函数取得最大值或最小值时自变量$x$的集合. 
    \begin{jieda}
    法一: 令$t=\tan{x}\in(-\infty,+\infty)$, 于是$y=1-\dfrac{2t}{t^2+t+1}$. 问题结束了. \\
    法二: $y=\dfrac{\sin^2x-\sin{x}\cos{x}+\cos^2x}{\sin^2x+\sin{x}\cos{x}+\cos^2x}=\dfrac{1-\frac{1}{2}\sin2x}{1+\frac{1}{2}\sin2x}\in[\frac{1}{3},3]$. \\
    $\min=\frac{1}{3}(\sin2x=1\iff x=k\pi+\frac{\pi}{4})$, $\max=3(\sin2x=-1\iff x=k\pi-\frac{\pi}{4})$. \\
    \ps 法二中定义域没考虑, 可以最后反过来检验一下. 
    \end{jieda}
\end{enumerate}
C组：
\begin{enumerate}
    \item 已知$A, B \in \left(0, \dfrac{\pi}{2}\right)$，且$\dfrac{sinA}{sinB} = sin(A+B)$，求$tanA$的最大值.
    \begin{jieda}
        $\dfrac{sinA}{sinB} = sin(A+B)$即$\dfrac{sinA}{sinB} = sinAcosB+cosAsinB$，即$sinA\left(\dfrac{1}{sinB} - cosB\right) = cosAsinB$，即$tanA = \dfrac{sinB}{\left(\dfrac{1}{sinB} - cosB\right)} = \dfrac{sin^2B}{1 - sinBcosB} = \dfrac{tan^2B}{tan^2B - tanB + 1}$，设$k = \dfrac{1}{tanB}$，则$\dfrac{tan^2B}{tan^2B - tanB + 1} = \dfrac{1}{1 - k + k^2} \leq \dfrac{4}{3}$，此时$tanB = 2$ \\
        \ps $\dfrac{sin^2B}{1 - sinBcosB} = \dfrac{1-cos2B}{2-sin2B}$，为$(1, 2)$与单位圆上半圆上点连线斜率的倒数，根据几何关系可知相切时可取最值，由正切和角公式可知最小的斜率为$\dfrac{3}{4}$，故$tanA$最大值为$\dfrac{4}{3}$，此时$sinB = \dfrac{2\sqrt{5}}{5}, cosB = \dfrac{\sqrt{5}}{5}$
    \end{jieda}
    \item 已知锐角$A, B$满足$tan(A+B) = 2tanA$，求$tanB$的最大值.
    \begin{jieda}
        $\dfrac{tanA + tanB}{1-tanA \cdot tanB} = 2tanA$，故$tanB = \dfrac{tanA}{1+2tan^2A} = \dfrac{1}{2tanA + \dfrac{1}{tanA}}$，因为$tanA > 0$，所以$tanB \leq \dfrac{\sqrt{2}}{4}$，当$tanA = \dfrac{\sqrt{2}}{2}$时取等号.
    \end{jieda}
    \item 在$\triangle ABC$中$tan\dfrac{A}{2} + tan\dfrac{B}{2} = 1$，求$tan\dfrac{C}{2}$的最小值.
    \begin{jieda}
        \begin{enumerate}
            \item 方法1 \\
            $tan\dfrac{C}{2} = cot\left(\dfrac{A+B}{2}\right) = \dfrac{1-tan\dfrac{A}{2}tan\dfrac{B}{2}}{tan\dfrac{A}{2}+tan\dfrac{B}{2}} = 1-tan\dfrac{A}{2}tan\dfrac{B}{2} \geq 1 - \left(\dfrac{tan\dfrac{A}{2} + tan\dfrac{B}{2}}{2}\right)^2 = \dfrac{3}{4}$，当$tan\dfrac{A}{2} = tan\dfrac{B}{2} = \dfrac{1}{2}$时取等.
            \item 方法2 \\
            $tan\dfrac{A}{2} + tan\dfrac{B}{2} = 1$即$\dfrac{sin\left(\dfrac{A+B}{2}\right)}{cos\dfrac{A}{2}cos\dfrac{B}{2}} = 1$，\\
            即$cos\dfrac{C}{2} = cos\dfrac{A}{2}\cdot cos\dfrac{B}{2} = \dfrac{1}{2}\left[cos\left(\dfrac{A+B}{2}\right) + cos\left(\dfrac{A-B}{2}\right)\right]$，即$2cos\dfrac{C}{2} = sin\dfrac{C}{2}+cos\left(\dfrac{A-B}{2}\right)$，即$cos\left(\dfrac{A-B}{2}\right) = 2cos\dfrac{C}{2} - sin\dfrac{C}{2} = 2\sqrt{\dfrac{1}{1+tan^2\dfrac{C}{2}}} - \sqrt{\dfrac{tan^2\dfrac{C}{2}}{1+tan^2\dfrac{C}{2}}} = \dfrac{2 - tan\dfrac{C}{2}}{\sqrt{1+tan^2\dfrac{C}{2}}}$，等式右侧单调递减，因此当左侧取得最大值1（$A = B$）时$tan\dfrac{C}{2}$取得最小值$\dfrac{3}{4}$ \\
            \ps 得到$cos\left(\dfrac{A-B}{2}\right) = 2cos\dfrac{C}{2} - sin\dfrac{C}{2}$时即可因右侧和$tan\dfrac{C}{2}$单调性相反判断出左侧取最大值时对应$tan\dfrac{C}{2}$的最小值.
        \end{enumerate}
    \end{jieda}
\end{enumerate}

\clearpage

\section{简单的三角方程}
\zsd 简单的三角方程

A组：
\begin{enumerate}
\item 写出下列$x$的解集，其中$a\in (-1,1)$
\begin{enumerate}[(1)]
    \item $\sin x=a$
    \item $\cos x=a$
    \item $\tan x =a$
\end{enumerate}
\begin{jieda}
    \begin{enumerate}
        \item $\{x | x = arcsin(a) + 2k\pi \mbox{或} x = \pi - arcsin(a) + 2k\pi, k \in \mathbf{Z}\}$
        \item $\{x | x = arccos(a) + 2k\pi \mbox{或} x = - arccos(a) + 2k\pi, k \in \mathbf{Z}\}$
        \item $\{x | x = arctan(a) + k\pi, k \in \mathbf{Z}\}$
    \end{enumerate}
\end{jieda}

\item $\sin\beta=\sin\alpha\iff$\blkda[7]{$\alpha = \beta + 2k\pi$ 或 $\alpha = \pi - \beta + 2k\pi, k \in \mathbf{Z}$}, \\
$\cos\beta=\cos\alpha\iff$\blkda[7]{$\alpha = \beta + 2k\pi$ 或 $\alpha = - \beta + 2k\pi, k \in \mathbf{Z}$}, \\
$\tan\beta=\tan\alpha\iff$\blkda[7]{$\alpha = \beta + k\pi, k \in \mathbf{Z}$}, \\

\end{enumerate}

B组：
\begin{enumerate}

\item $8\sin^2x+3\sin2x-4=0$的解集为\blkda[8]{$\{x\mid x=k\pi-\arctan2\text{或}k\pi+\arctan\frac{1}{2}, k\zz\}$}. 
\begin{jieda}
$8\sin^2x+6\sin{x}\cos{x}-4\sin^2x-4\cos^2x=0\\
\iff 2\tan^2x+3\tan{x}-2=0\iff\tan{x}=-2$或$\frac{1}{2}$. \\
亦可以: $3\sin2x-4\cos2x=0\iff 5\sin(2x-\arccos\frac{3}{5})=0$. \\
也可以: $3\sin2x-4\cos2x=0\iff \tan2x=\frac{4}{3}\iff x=\frac{k\pi}{2}+\frac{1}{2}\arctan\frac{4}{3}$. 
\end{jieda}

\item $\sqrt{1-\sin^2x}=\sin{x}$的解集为\blkda[6]{$\{x\mid x=2k\pi+\frac{\pi}{4}\text{或}2k\pi+\frac{3\pi}{4}, k\zz\}$}. 
\begin{jieda}
$\sin{x}=\frac{\sqrt{2}}{2}$.
\end{jieda}

\item 方程$\sqrt{3}\sin2x-\cos2x+1=0$在$[0,2\pi)$上的解集为\blkda[3]{$\{0,\pi,\frac{2\pi}{3},\frac{5\pi}{3}\}$}. 
\begin{jieda}
角度一: $2\sin(2x-\frac{\pi}{6})+1=0$. \\
角度二: 万能公式. \\
角度三: 结合$\sin^22x+\cos^22x=1$, 解得$(\cos2x,\sin2x)=(1,0),(-\frac{1}{2},-\frac{\sqrt{3}}{2})$. 
\end{jieda}


\item 在$[0,\frac{3\pi}{2}]$内, 方程$\cos(\pi\cos{x})=0$的所有解的和为\blkda{$\frac{7\pi}{3}$}. 
\begin{jieda}
$\cos(\pi\cos{x})=0\iff\pi\cos{x}=k\pi+\frac{\pi}{2}\iff\cos{x}=k+\frac{1}{2}$. \\
又$\cos{x}\in[-1,1]$, 故$k=-1,0$, 即$\cos{x}=\pm\frac{1}{2}$. \\
又$x\in[0,\frac{3\pi}{2}]$, 故$x=\frac{\pi}{3},\frac{2\pi}{3},\frac{4\pi}{3}$. 
\end{jieda}

\item 方程$\lg(\sqrt{3}\sin{x})=\lg(-\cos{x})$的解集为\blkda[4]{$\{x\mid x=2k\pi+\frac{5\pi}{6},k\zz\}$}. 
\begin{jieda}
$\sqrt{3}\sin{x}=-\cos{x}>0$, 故$x=2k\pi+\frac{5\pi}{6}$. 
\end{jieda}

\item $\tan^2x-3\tan{x}+2=0$的解集为\blkda[7]{$\{x\mid x=k\pi+\frac{\pi}{4}\text{或}k\pi+\arctan2,k\zz\}$}. 
\begin{jieda}
$\iff\tan{x}=1,2$. 
\end{jieda}

\item 方程$\sin{x}=\frac{1}{2}(a+\frac{1}{a})$($a\zr$且$a\neq0)$在$[0,2\pi)$内\dotfill(\kh{A})
\xx{可能有一解}{可能有两解}{有无穷多解}{一定无解}
\begin{jieda}
$\frac{1}{2}(a+\frac{1}{a})\in(-\infty,-1]\cup[1,+\infty)$, $\sin{x}\in[-1,1]$, \\
于是若有解, 就只能$\sin{x}=1$或$-1$(只能取其一). 
\end{jieda}

\item 已知$2\sin{x}-2\sqrt{3}\cos{x}=\dfrac{6-2m}{3+m}$, 则$m$的取值范围是\dotfill(\kh{D})
\xx{$[-9,-1]$}
{$[-9,+\infty)$}
{$[-1,+\infty)$}
{$(-\infty,-9]\cup[-1,+\infty)$}
\begin{jieda}
$\dfrac{6-2m}{3+m}=2\sin{x}-2\sqrt{3}\cos{x}=4\sin(x-\frac{\pi}{3})\in[-2,2]$. 
\end{jieda}

\item 三角方程$\tan3x+\tan2x=0$在$[0,\pi]$内解的个数为\dotfill(\kh{A})
\xx{6}{5}{4}{3}
\begin{jieda}
$\tan3x+\tan2x=0\iff \tan3x=\tan(-2x)$. \\
于是$3x=-2x+k\pi$或$3x=2x+k\pi$, 故$x=\frac{k\pi}{5}$或$k\pi$. \\
代入检验定义域, 可知解为$0,\pi,\frac{\pi}{5},\frac{2\pi}{5},\frac{3\pi}{5},\frac{4\pi}{5}$. 
\end{jieda}

\item 已知$y=f(x)$是周期为$2\pi$的函数, 当$x\in[0,2\pi)$时, $f(x)=\sin\frac{x}{2}$, 则$f(x)=\frac{1}{2}$的解集为\dotfill(\kh{CD})
\xx
{$\left\{x\mid x=2k\pi+\frac{\pi}{3},k\zz\right\}$}
{$\left\{x\mid x=2k\pi+\frac{5\pi}{3},k\zz\right\}$}
{$\left\{x\mid x=2k\pi\pm\frac{\pi}{3},k\zz\right\}$}
{$\left\{x\mid x=2k\pi\pm\frac{5\pi}{3},k\zz\right\}$}
\begin{jieda}
在$[0,2\pi)$上, 解$\sin\frac{x}{2}=\frac{1}{2}$得$x=\frac{\pi}{3}$或$\frac{5\pi}{3}$. 
结合周期, 可得CD等效. 
\end{jieda}

\item 已知关于$x$的方程$x^2-(\sin\alpha+\cos\alpha)x+\sin^2\alpha-\sin\alpha\cos\alpha-1=0$有两个相等的实根, 求$\alpha$与$x$的值. 
\begin{jieda}
$\Delta=(\sin\alpha+\cos\alpha)^2-4(\sin^2\alpha-\sin\alpha\cos\alpha-1)=\sin^2\alpha+6\sin\alpha\cos\alpha+5\cos^2\alpha=0$, 
故$\tan\alpha=-1$或$-5$, 于是$\alpha=k\pi-\frac{\pi}{4}$或$k\pi-\arctan5(k\zz)$. \\
当$\alpha=k\pi-\frac{\pi}{4}$时, 求得$x=0$; $\alpha=k\pi-\arctan5$时, 求得$x=\pm\frac{\sqrt{26}}{13}$. 
\end{jieda}

\item 已知$\alpha,\beta$是方程$\sin{x}+\sqrt{3}\cos{x}=m$在$(0,2\pi)$上的两相异实根. 
\begin{enumerate}[(1)]
\item 求$m$的取值范围; 
\item 求$\alpha+\beta$的值.
\end{enumerate}
\begin{jieda}
$\iff\sin(x+\frac{\pi}{3})=\frac{m}{2}$. 作$y=\sin(x+\frac{\pi}{3})$在$(0,2\pi)$上的草图如下: \\
\begin{minipage}{\linewidth}
\centering
\includegraphics[totalheight=6\baselineskip]{./bodyxb/src/main42xbb2ch07sanjiaofangchengt13.png}
\end{minipage}
\begin{enumerate}[(1)]
\item $\frac{m}{2}\in(-1,\frac{\sqrt{3}}{2})\cup(\frac{\sqrt{3}}{2},1)$, 故$m\in(-2,\sqrt{3})\cup(\sqrt{3},2)$. 
\item 当$m\in(\sqrt{3},2)$时, $\alpha+\beta=\frac{\pi}{3}$; 
当$m\in(-2,\sqrt{3})$时, $\alpha+\beta=\frac{7\pi}{3}$. 
\end{enumerate}
\end{jieda}

\item 若关于$x$的方程$2\cos^2(x+\pi)-\sin{x}+a=0$有实根, 求实数$a$的取值范围. 
\begin{jieda}
$\iff a=-2\cos^2x+\sin{x}=2\sin^2x+\sin{x}-2=2(\sin{x}+\frac{1}{4})^2-\frac{17}{8}\in[-\frac{17}{8},1]$. 
\end{jieda}

\item 如图，一智能扫地机器人在 A 处发现位于它正西方向的 $B$ 处和北偏东 $\frac{\pi}{6}$ 方向上的 $C$ 处分别有需要清扫的垃圾，红外线感应测量发现机器人到 $B$ 的距离比到 $C$的距离少 0.4 米，于是选择沿 $A \rightarrow B \rightarrow C$ 路线清扫，已知智能扫地机器人的直线行走速度为 0.2 米／秒， 10 秒钟完成了清扫任务（忽略机器人吸入垃圾及在 $B$ 处旋转所用时间）

\begin{minipage}{\linewidth}
\centering
\includegraphics[totalheight=6\baselineskip]{./bodyxb/src/zhaopianch6jiandande.png}
\end{minipage}

\begin{enumerate}[(1)]
\item $B,C$ 两处垃圾的距离是多少？
\item 求智能扫地机器人此次清扫过程中旋转角的最小值？请指出旋转方向．

\begin{jieda}
    \begin{enumerate}[(1)]
\item 依题意可知 $\angle A=\frac{\pi}{2}+\frac{\pi}{6}=\frac{2 \pi}{3}$ ，设 $A, B, C$ 的对边分别为 $a, b, c$ ，则
$$
\left\{\begin{array} { l } 
{ a + c = 0 . 2 \times 1 0 } \\
{ b - c = 0 . 4 } \\
{ a ^ { 2 } = b ^ { 2 } + c ^ { 2 } - 2 b c \operatorname { c o s } A }
\end{array} \Rightarrow \left\{\begin{array}{l}
a+c=2 \\
b-c=0.4 \\
a^2=b^2+c^2+b c
\end{array},\right.\right.
$$
将 $c=2-a, b=c+0.4=2.4-a$ 代人 $a^2=b^2+c^2+b c$ 化简得 $a^2-6.6 a+7.28=0$ ，即 $(a-1.4)(a-5.2)=0$ ，

由于 $a<0.2 \times 10=2$ ，所以 $(a-1.4)(a-5.2)=0$ 解得 $a=1.4 \mathrm{~m}$ ．
且 $c=2-a=0.6 \mathrm{~m}, b=2.4-a=1 \mathrm{~m}$ ．
\item 由余弦定理得 $\cos B=\frac{a^2+c^2-b^2}{2 a c}=\frac{1.96+0.36-1}{2 \times 1.4 \times 0.6}=\frac{1.32}{1.68}=\frac{11}{14}$ ，由于 $A=\frac{2 \pi}{3}$ ，所以 $B$ 为锐角．
所以旋转角的最小值为 $\pi-\arccos \frac{11}{14}$ ，方向沿顺时针方向旋转．
    \end{enumerate}
\end{jieda}
\end{enumerate}

C组：
\begin{enumerate}
    \item 已知 $a, b, c$ 分别为 $\triangle A B C$ 三个内角 $A, B, C$ 的对边， $\mathrm{b} \cos C+\sqrt{3} b \sin C-a-c=0$ ．
    \begin{enumerate}[(1)]
        \item 求证 $A, B, C$ 成等差数列；
        \item 若 $b=\sqrt{3}$ ，求 $2 a+c$ 的最大值．
    \end{enumerate}
\begin{jieda}
    \begin{enumerate}[(1)]
        \item 因为 $\mathrm{b} \cos C+\sqrt{3} b-a-c=0$ ，所以
$$
\begin{aligned}
& \sin B \cos C+\sqrt{3} \sin B \sin C-\sin A-\sin C=0 . \\
& \text { 又 } A+B+C=\pi \text {, 所以 } \sin A=\sin (B+C) \text {.所以 } \\
& \sin B \cos C+\sqrt{3} \sin B \sin C-\sin (B+C)-\sin C=0 \\
& \sin B \cos C+\sqrt{3} \sin B \sin C-(\sin B \cos C+\cos B \sin C)-\sin C=0
\end{aligned}
$$

$$
\sqrt{3} \sin B \sin C-\cos B \sin C-\sin C=0
$$

又 $0<C<\pi$ ，所以 $\sin C>0$ ．所以

$$
\sqrt{3} \sin B-\cos B-1=0
$$

$$
2\left(\sin B \cdot \frac{\sqrt{3}}{2}-\cos B \cdot \frac{1}{2}\right)=1
$$

$$
2\left(\sin B \cos \frac{\pi}{6}-\cos B \sin \frac{\pi}{6}\right)=1 \text { 亦即 } 2\sin \left(B-\frac{\pi}{6}\right)=1
$$
又 $0<B<\pi$ ，所以 $B=\frac{\pi}{3}$ ．
又 $\mathrm{A}+\mathrm{B}+\mathrm{C}=\pi$ ，所以 $\mathrm{A}+\mathrm{C}=\frac{2 \pi}{3}=2 B$ ．
所以 $A, ~ B, ~ C$ 成等差数列．

\item 由正弦定理 $\frac{a}{\sin A}=\frac{b}{\sin B}=\frac{c}{\sin C}$ ，得

$$
\frac{a}{\sin A}=\frac{c}{\sin \bar{C}}=\frac{\sqrt{3}}{\frac{\sqrt{3}}{2}}
$$


即 $a=2 \sin A, \mathrm{c} \sin C$ ．所以 $2 a+c=4 \sin A+2 \sin C$ ．
因为 $\mathrm{B}=\frac{\pi}{3}$ ，所以 $\mathrm{C}=\frac{2 \pi}{3}-A$ ．所以

$$
\begin{aligned}
2 a+c&=4 \sin A+2 \sin \left(\frac{2 \pi}{3}-A\right) \\
& =4 \sin A+2\left(\frac{\sqrt{3}}{2} \cos A+\frac{1}{2} \sin A\right) \\
& =5 \sin A+\sqrt{3} \cos A \\
&=2 \sqrt{7}\left(\sin A \cdot \frac{5 \sqrt{7}}{14}+\cos A \cdot \frac{\sqrt{21}}{14}\right) \\
& =2 \sqrt{7} \sin (A+\varphi) \leq 2 \sqrt{7}
\end{aligned}
$$


其中锐角 $\varphi$ 满足 $\sin \varphi=\frac{\sqrt{21}}{14}, \cos \varphi=\frac{5 \sqrt{7}}{14}$ ．

因为 $A \in\left(0, \frac{2 \pi}{3}\right)$ ，所以当 $A$ 为 $\varphi$ 的余角时，即 $2 \sqrt{7} \sin (A+\emptyset)=2 \sqrt{7}$ 时， $2 a+c$ 取最大值 $2 \sqrt{7}$ ．
    \end{enumerate}
\end{jieda}
\end{enumerate}\end{enumerate}